\chapter{Mathematical Tools}\label{ch:tools}

Quantum mechanics is written in the language of linear algebra, and a quantum
computer is, in the end, a machine for doing linear algebra on a vector that
nobody is allowed to look at. This chapter builds that language from the
ground up, with the quantum applications in view from the start.

We begin with the physical puzzles that forced the theory upon us, and show
that the one equation that resolved them, Schr\"odinger's, is \emph{linear}
(Section~\ref{sec:motivation}). Everything else follows from taking that word
seriously. We review the complex numbers, which are the scalars of the theory
(Section~\ref{sec:complex}); vector spaces and linear maps, with matrices as
their coordinates (Section~\ref{sec:vector-spaces}); inner products, which
give lengths, angles and probabilities, together with Dirac's notation
(Section~\ref{sec:hilbert}); and finally operators, whose eigenvalues will be
the outcomes of measurements and whose exponentials will be the gates of a
quantum computer (Section~\ref{sec:operators}).

\begin{objectives}[In this chapter]
  You will compute with complex numbers in Cartesian and polar
  form; decide whether a set is a vector space, find a basis of it, and write a
  linear map as a matrix; compute inner products, norms and orthonormal bases
  in Dirac notation; and diagonalise Hermitian and unitary operators and
  apply functions, such as exponentials, to them.
\end{objectives}

\begin{whereitsits}
  This is the first half of Part~II. It is self-contained and assumes only
  basic calculus and some familiarity with matrices. The second half,
  Chapter~\ref{ch:postulates}, uses every tool built here to state the four
  postulates of quantum mechanics.
\end{whereitsits}

%======================================================================
\section{Why quantum?}\label{sec:motivation}
%======================================================================

Around 1900, physics looked nearly finished. Newton's mechanics described
planets and projectiles, Maxwell's equations described light and electricity,
and thermodynamics described heat. A handful of experiments refused to fit.
Three of them, taken together, are enough to see why a new theory was needed,
and what kind of mathematics it had to use.

\subsection{Three experiments}

\paragraph{The double slit.}
Send a stream of particles towards a wall with two narrow slits, and record
where they land on a screen behind it. If the particles were little bullets,
each one would go through one slit or the other, and the pattern on the screen
would be the sum of two bumps, one behind each slit. Waves behave differently.
A wave goes through \emph{both} slits, the two outgoing waves overlap, and the
screen shows alternating bright and dark bands, called \emph{interference
fringes}, because the waves reinforce each other at some points and cancel at
others (Figure~\ref{fig:slit}).

\begin{figure}[htb]
  \centering
  \begin{minipage}[c]{0.48\textwidth}
    \centering
    \begin{tikzpicture}[scale=0.82]
      % barrier with two slits at y = +-0.5
      \draw[ultra thick,ink] (0,-2.3)--(0,-0.6);
      \draw[ultra thick,ink] (0,-0.4)--(0,0.4);
      \draw[ultra thick,ink] (0,0.6)--(0,2.3);
      % source and incoming rays
      \fill[navy] (-2,0) circle (2.5pt) node[below=3pt,font=\footnotesize] {source};
      \draw[navy,->] (-2,0)--(-0.1,0.5);
      \draw[navy,->] (-2,0)--(-0.1,-0.5);
      % outgoing circular waves from each slit
      \foreach \r in {0.6,1.1,1.6,2.1,2.6,3.1,3.6}{
        \draw[navy!45] ($(0,0.5)+(-50:\r)$) arc (-50:50:\r);
        \draw[navy!45] ($(0,-0.5)+(-50:\r)$) arc (-50:50:\r);
      }
      % screen with fringes drawn as shades
      \foreach \y in {-2.2,-2.15,...,2.15}{
        \pgfmathsetmacro{\o}{0.5*(1+cos(\y*360/1.1))}
        \fill[navy,opacity=\o] (4.05,\y) rectangle (4.5,{\y+0.05});
      }
      \draw[ultra thick,ink] (4,-2.3)--(4,2.3);
      \node[font=\footnotesize] at (4.3,-2.65) {screen};
    \end{tikzpicture}
  \end{minipage}\hfill
  \begin{minipage}[c]{0.48\textwidth}
    \centering
    \begin{tikzpicture}
      \begin{axis}[width=\linewidth, height=5.2cm, xmin=-2.2, xmax=2.2, ymin=0, ymax=4.6,
          xlabel={position $x$ on the screen}, ylabel={intensity},
          ytick=\empty, xtick=\empty, axis lines=left, tick style={draw=none},
          legend style={font=\footnotesize, at={(0.5,1.02)}, anchor=south, draw=none},
          label style={font=\footnotesize}]
        \addplot[navy, very thick, samples=200, domain=-2.2:2.2] {2*(1+cos(deg(2*pi*x/1.1)))};
        \addplot[figacc, very thick, dashed, domain=-2.2:2.2] {2};
        \legend{waves: $\abs{\psi_1+\psi_2}^2$, bullets: $\abs{\psi_1}^2+\abs{\psi_2}^2$}
      \end{axis}
    \end{tikzpicture}
  \end{minipage}
  \caption{Left: the double-slit set-up, with the fringes drawn on the screen. Right: the
  pattern that waves produce (solid) against the pattern that bullets would produce (dashed).}
  \label{fig:slit}
\end{figure}

The surprise is that electrons, photons and even whole molecules produce the
\emph{fringes}, and they do so even when they are sent through the apparatus
one at a time. Each particle arrives at a single point, like a bullet, but
after thousands of arrivals the points build up the interference pattern, as if
each particle had passed through both slits and interfered with itself.

Mathematically this is easy to express if we give each route a complex number.
Let $\psi_1(x)=a\,e^{\ii\varphi_1(x)}$ and $\psi_2(x)=a\,e^{\ii\varphi_2(x)}$ be
the \term{amplitudes} for reaching the point $x$ through slit~1 and slit~2. The
rule is that amplitudes \emph{add}, and the intensity is the squared modulus of
the total:
\[
  \abs{\psi_1+\psi_2}^2=2a^2\bigl(1+\cos(\varphi_1-\varphi_2)\bigr),
\]
which oscillates between $0$ and $4a^2$ as the phase difference changes. This
is precisely the computation of Example~\ref{ex:relative-phase} in the next section: complex numbers
and their phases are where the fringes come from. Bullets would give
$\abs{\psi_1}^2+\abs{\psi_2}^2=2a^2$, the flat dashed line.

\paragraph{Black-body radiation.}
A hot object glows: first red, then orange, then white. A perfect absorber
and emitter of radiation, called a \emph{black body}, glows with a
characteristic spectrum that depends only on its temperature. Classical physics
predicts that the energy emitted at frequency $\nu$ grows like $\nu^2$, so the
total energy would be infinite, the so-called \emph{ultraviolet catastrophe}.
Real bodies do not do this: their curves rise, peak, and fall again.

In 1900 Planck found a formula that matches the measurements exactly, and the
price of the formula was a strange hypothesis: light of frequency $\nu$ can
only be exchanged with matter in whole units of energy $E=h\nu$. With this
assumption the energy at frequency $\nu$ is proportional to
$\nu^3/(e^{h\nu/k_BT}-1)$, which agrees with the classical $\nu^2$ law at low
frequencies and is cut off at high ones (Figure~\ref{fig:blackbody}).

\begin{figure}[htb]
  \centering
  \begin{tikzpicture}
    \begin{axis}[width=0.8\linewidth, height=5.4cm, xmin=0, xmax=8, ymin=0, ymax=3,
        xlabel={frequency, in units of $k_BT/h$}, ylabel={emitted energy},
        ytick=\empty, xtick={0,2,4,6,8}, axis lines=left,
        legend style={font=\footnotesize, at={(0.04,0.96)}, anchor=north west, draw=none},
        label style={font=\footnotesize}, tick label style={font=\footnotesize}]
      \addplot[figacc, very thick, dashed, domain=0:1.74, samples=60] {x^2};
      \addplot[navy, very thick, samples=120, domain=0.01:8] {x^3/(exp(x)-1)};
      \legend{classical, Planck}
    \end{axis}
  \end{tikzpicture}
  \caption{Spectrum of a black body. The classical prediction (dashed, cut off at the plot's top)
  keeps growing; Planck's formula (solid) peaks and falls. The curves agree at low frequency.}
  \label{fig:blackbody}
\end{figure}

\paragraph{The stability of the atom.}
In Rutherford's picture of the hydrogen atom, an electron orbits a nucleus. But
an orbiting charge is an accelerating charge, and Maxwell's equations say that
accelerating charges radiate. Classically the electron would lose energy and
spiral into the nucleus in about $10^{-11}$ seconds: atoms could not exist.
Moreover, hydrogen gas does not glow with a continuous spectrum but with a
handful of sharp lines, at frequencies that fit the pattern
$E_n=-13.6\,\mathrm{eV}/n^2$ for the energy levels of the atom
(Figure~\ref{fig:levels}): only certain energies are allowed, and light is
emitted exactly when the electron jumps between them.

\begin{figure}[htb]
  \centering
  \begin{tikzpicture}[scale=1.0]
    % vertical positions are not to scale: n=1 -> -4.6, n=2 -> -1.6, n=3 -> -0.95, n=4 -> -0.55
    \foreach \n/\E/\y in {1/-13.6/-4.6,2/-3.4/-1.6,3/-1.51/-0.95,4/-0.85/-0.55}{
      \draw[ink, thick] (0,\y) -- (4.2,\y);
      \node[left,font=\footnotesize] at (0,\y) {$n=\n$};
      \node[right,font=\footnotesize] at (4.2,\y) {$\E$\,eV};
    }
    \draw[ink, thick, dotted] (0,0) -- (4.2,0);
    \node[right,font=\footnotesize] at (4.2,0) {$0$};
    \draw[->,navy,thick] (0.7,-1.6) -- (0.7,-4.6);
    \draw[->,navy,thick] (1.3,-0.95) -- (1.3,-4.6);
    \draw[->,navy,thick] (1.9,-0.55) -- (1.9,-4.6);
    \node[font=\footnotesize,navy,align=center] at (1.3,-5.2) {ultraviolet lines\\(towards $n=1$)};
    \draw[->,figacc,thick] (2.9,-0.95) -- (2.9,-1.6);
    \draw[->,figacc,thick] (3.5,-0.55) -- (3.5,-1.6);
    \node[font=\footnotesize,figacc,align=center] at (3.2,-2.4) {visible lines\\(towards $n=2$)};
    \node[font=\scriptsize,ink] at (2.1,0.35) {(heights not to scale)};
  \end{tikzpicture}
  \caption{Energy levels of hydrogen. Only the energies shown are allowed. Light is emitted
  when the electron drops from one level to a lower one; the drops towards $n=2$ give the
  visible lines.}
  \label{fig:levels}
\end{figure}

Each experiment has the same moral. Energy comes in units, behaviour is
described by amplitudes that add, and what we observe are squared moduli of
those amplitudes. We need a theory whose basic objects are complex-valued
\emph{waves of probability amplitude}, and whose basic rule is addition.

\subsection{Schr\"odinger's equation}

In 1926 Schr\"odinger wrote down the equation governing such waves for a
particle of mass $m$ moving in a potential $V$:
\begin{equation}\label{eq:schrodinger}
  \ii\hbar\,\frac{\partial\psi}{\partial t}(x,t)
  =\hat H\psi(x,t),\qquad
  \hat H=-\frac{\hbar^2}{2m}\frac{\partial^2}{\partial x^2}+V(x).
\end{equation}
The unknown $\psi(x,t)$ is complex-valued, and it is called the \term{wave
function}. According to Born, $\abs{\psi(x,t)}^2$ is the probability density of
finding the particle near $x$ at time $t$. We take this as motivation here; in
Chapter~\ref{ch:postulates} it becomes a postulate.

Let us look at the simplest case, a free particle ($V=0$).

\begin{example}\label{ex:planewave}
  Show that $\psi(x,t)=A\,e^{\ii(kx-\omega t)}$ with $A\in\C$ solves the free
  equation if $\hbar\omega=\hbar^2k^2/2m$.
\end{example}
\begin{solution}
  Differentiate: $\partial_t\psi=-\ii\omega\psi$ and $\partial_x^2\psi=-k^2\psi$. Substituting
  into $\ii\hbar\,\partial_t\psi=-\frac{\hbar^2}{2m}\partial_x^2\psi$,
  \[
    \ii\hbar(-\ii\omega)\psi=\hbar\omega\,\psi
    \qquad\text{and}\qquad
    -\frac{\hbar^2}{2m}(-k^2)\psi=\frac{\hbar^2k^2}{2m}\,\psi,
  \]
  so the equation holds exactly when $\hbar\omega=\hbar^2k^2/2m$. Note also that
  $\abs{\psi}^2=\abs{A}^2$ does not depend on $x$ or $t$.
\end{solution}

The relation between $\omega$ and $k$ is the free particle's energy
$E=\hbar\omega=p^2/2m$ with momentum $p=\hbar k$. That a plane wave has a
constant $\abs{\psi}^2$ also reveals a difficulty: it cannot be normalised
over all of $\R$, so it is an idealisation.

Now the observation that drives this whole chapter. Equation
\eqref{eq:schrodinger} is \emph{linear} in $\psi$.

\begin{example}\label{ex:linear-schrodinger}
  Show that if $\psi_1$ and $\psi_2$ solve \eqref{eq:schrodinger}, so does
  $\alpha\psi_1+\beta\psi_2$ for all $\alpha,\beta\in\C$. Does the product $\psi_1\psi_2$ solve it?
\end{example}
\begin{solution}
  Both $\partial_t$ and $\hat H$ are linear operations, so
  \[
    \ii\hbar\,\partial_t(\alpha\psi_1+\beta\psi_2)
    =\alpha\,\ii\hbar\,\partial_t\psi_1+\beta\,\ii\hbar\,\partial_t\psi_2
    =\alpha\hat H\psi_1+\beta\hat H\psi_2
    =\hat H(\alpha\psi_1+\beta\psi_2).
  \]
  The product does not work in general: $\hat H$ contains a second derivative,
  and $(\psi_1\psi_2)''=\psi_1''\psi_2+2\psi_1'\psi_2'+\psi_1\psi_2''$ produces a cross term
  that has no counterpart on the other side.
\end{solution}

So the set of solutions is closed under sums and complex multiples: it is a
\emph{complex vector space}. That is the same statement as the double slit, in
which the amplitude for ``both slits'' is the \emph{sum} of the amplitudes for
each. Superposition is not a mysterious extra; it is linearity.

Finally, one special family of solutions will matter a great deal. If $\hat H$
does not depend on time, look for solutions of the form
$\psi(x,t)=\varphi(x)\,T(t)$. Substituting, one finds that $\varphi$ must satisfy
$\hat H\varphi=E\varphi$ for a constant $E$, and that $T(t)=e^{-\ii Et/\hbar}$
(Exercise~\ref{ex:separation}). Thus
\[
  \psi(x,t)=\varphi(x)\,e^{-\ii Et/\hbar},\qquad \hat H\varphi=E\varphi .
\]
The equation $\hat H\varphi=E\varphi$ is an \emph{eigenvalue problem}: $\varphi$ is an eigenvector of
the operator $\hat H$, and $E$, the energy, is its eigenvalue. The discrete
levels of Figure~\ref{fig:levels} are exactly the eigenvalues of the atom's
Hamiltonian. These \term{stationary states} only get multiplied by a phase as
time goes on, so $\abs{\psi}^2$ does not change.

This is the programme of the chapter. Solutions of a linear equation live in
a vector space; probabilities come from complex numbers through their moduli;
and the questions we can ask about the system are eigenvalue problems. Each of
these three ideas gets its own section. We begin, in the next one, with the
numbers.

\begin{exerciseset}
  \exercise{1} In the double slit, let $\psi_j(x)=a\,e^{\ii\varphi_j(x)}$ with $a>0$. Show that
  $\abs{\psi_1+\psi_2}^2=2a^2\bigl(1+\cos(\varphi_1-\varphi_2)\bigr)$, and find the phase
  differences for which the screen is dark.

  \exercise{1} Let $\psi=A\,e^{\ii(kx-\omega t)}$. Compute $\partial_t\psi$ and $\partial_x^2\psi$,
  substitute into the free equation, and derive $\omega=\hbar k^2/2m$. Is $\psi$ normalisable over $\R$?

  \exercise{2} Let $\hat H=-\frac{\hbar^2}{2m}\partial_x^2+V(x)$. Prove that the set of solutions
  of $\ii\hbar\,\partial_t\psi=\hat H\psi$ is a complex vector space. Give a concrete example of
  two solutions of the free equation whose product is not a solution.

  \exercise{2}\label{ex:separation} Assume $\hat H$ does not depend on $t$ and look for solutions
  $\psi(x,t)=\varphi(x)T(t)$. Show that $\ii\hbar\,T'(t)/T(t)=(\hat H\varphi)(x)/\varphi(x)$, so both
  sides equal a constant $E$. Solve for $T$, and prove that $\abs{\psi}^2$ does not depend on $t$.

  \exercise{2} Let $\hat H\varphi_j=E_j\varphi_j$ for $j=1,2$ with $E_1\neq E_2$ and $\varphi_j$
  real-valued, and let $\psi=c_1\varphi_1e^{-\ii E_1t/\hbar}+c_2\varphi_2e^{-\ii E_2t/\hbar}$ with
  $c_1,c_2>0$. Show that $\psi$ solves the time-dependent equation, and that $\abs{\psi}^2$
  oscillates with angular frequency $(E_2-E_1)/\hbar$. Is $\psi$ a stationary state?
\end{exerciseset}

%======================================================================
\section{Complex numbers}\label{sec:complex}
%======================================================================

\begin{objectives}
  You will add, multiply, conjugate and invert complex numbers, in Cartesian and
  in polar form; use the modulus and the triangle inequality; find the argument of a
  complex number without picking the wrong angle; and recognise a \emph{phase}, the
  idea that will separate observable from unobservable in quantum mechanics.
\end{objectives}

The state of a quantum system is a list of complex numbers called
\emph{amplitudes}. Probabilities will be their squared moduli, and the
difference between a ``global'' and a ``relative'' phase
(Section~\ref{sec:measurement}) will be the difference between what we can and
cannot observe. We review here the few facts about $\C$ that we will use all the time.

\subsection{Definition and arithmetic}

Some polynomials, such as $x^2+1$, have no root in $\R$. We add one by
decree and see what we get.

\begin{definition}[Complex numbers]\label{def:complex}
  The \term{imaginary unit} $\ii$ is a symbol with $\ii^2=-1$. A
  \term{complex number} is an expression $z=a+b\ii$ with $a,b\in\R$; we call
  $a=\Re z$ its \term{real part} and $b=\Im z$ its \term{imaginary part}. The set of
  all complex numbers is $\C$. Sums and products are
  \begin{align*}
    (a+b\ii)+(c+d\ii) &= (a+c)+(b+d)\ii,\\
    (a+b\ii)(c+d\ii) &= (ac-bd)+(ad+bc)\ii .
  \end{align*}
\end{definition}

The product is just ``expand the brackets and use $\ii^2=-1$''. We identify
$a\in\R$ with $a+0\ii$, so that $\R\subset\C$. A complex number is also a point
$(a,b)$ of the plane, which is how we draw it (Figure~\ref{fig:plane}).

\begin{figure}[htb]
  \centering
  \begin{tikzpicture}[>=stealth, scale=1.1]
    \draw[->,ink] (-1.3,0) -- (4.3,0) node[right] {$\Re$};
    \draw[->,ink] (0,-2.3) -- (0,2.6) node[above] {$\Im$};
    \coordinate (z) at (3,2);
    \coordinate (zb) at (3,-2);
    \draw[very thick, navy, ->] (0,0) -- (z);
    \draw[very thick, figacc, ->] (0,0) -- (zb);
    \draw[dashed, ink!60] (zb) -- (3,2);
    \fill[navy] (z) circle (1.8pt) node[above right] {$z=a+b\ii$};
    \fill[figacc] (zb) circle (1.8pt) node[below right] {$\bar z=a-b\ii$};
    \draw[ink!60] (3,0.08) -- (3,-0.08) node[below left, font=\small] {$a$};
    \draw[ink!60] (0.08,2) -- (-0.08,2) node[left, font=\small] {$b$};
    \draw[navy] (0.8,0) arc (0:33.7:0.8) node[midway, right, font=\small] {$\theta$};
    \node[navy, font=\small, rotate=33.7] at (1.35,1.2) {$\abs{z}$};
  \end{tikzpicture}
  \caption{A complex number $z=a+b\ii$ as a point of the plane, its conjugate $\bar z$ (the
  reflection in the real axis), its modulus $\abs{z}$ (the length of the arrow) and its argument $\theta$.}
  \label{fig:plane}
\end{figure}

Two small computations show the flavour. To multiply, expand and replace
$\ii^2$ by $-1$:
\[
  (2+\ii)(3-4\ii)=6-8\ii+3\ii-4\ii^2=10-5\ii .
\]
To raise to a power, it can be easier to square repeatedly: $(1+\ii)^2=2\ii$, so
$(1+\ii)^4=(2\ii)^2=-4$. The powers of $\ii$ itself repeat with period four,
$\ii^0=1,\ \ii^1=\ii,\ \ii^2=-1,\ \ii^3=-\ii,\ \ii^4=1,\dots$, which makes
large powers easy.

\begin{example}
  Compute $\ii^{2026}$.
\end{example}
\begin{solution}
  Only the remainder of the exponent modulo $4$ matters. Since $2026=4\cdot506+2$,
  $\ii^{2026}=\ii^2=-1$.
\end{solution}

Complex numbers obey the same algebraic rules as the reals: addition and
multiplication are associative and commutative, multiplication distributes over
addition, $0$ and $1$ are neutral, every $z=a+b\ii$ has an additive inverse
$-z=-a-b\ii$, and every $z\neq0$ has a multiplicative inverse. In the language
of algebra, $\C$ is a \emph{field}. All of this is a routine check from
Definition~\ref{def:complex}, except the inverse, which needs a formula; the
next subsection gives it.

\subsection{Conjugate and modulus}

\begin{definition}[Conjugate and modulus]\label{def:conj}
  For $z=a+b\ii$, its \term{conjugate} is $\bar z=a-b\ii$ and its
  \term{modulus} is $\abs{z}=\sqrt{a^2+b^2}$.
\end{definition}

Geometrically, $\bar z$ is the reflection of $z$ in the real axis and $\abs{z}$
is the distance from $z$ to the origin. The two are tied together by the following facts.

\begin{theorem}[Basic properties]\label{thm:basic}
  For all $z,w\in\C$:
  \begin{enumerate}
    \item[(a)] $\overline{z+w}=\bar z+\bar w$ and $\overline{zw}=\bar z\,\bar w$;
    \item[(b)] $z\bar z=\abs{z}^2$;
    \item[(c)] $\abs{zw}=\abs{z}\,\abs{w}$;
    \item[(d)] $\abs{z}=0$ if and only if $z=0$;
    \item[(e)] if $z\ne0$, then $z^{-1}=\bar z/\abs{z}^2$.
  \end{enumerate}
\end{theorem}
\begin{proof}
  (a) The statement about sums is immediate. For products let $z=a+b\ii$, $w=c+d\ii$. Then
  $zw=(ac-bd)+(ad+bc)\ii$, so $\overline{zw}=(ac-bd)-(ad+bc)\ii$, while
  \[
    \bar z\,\bar w=(a-b\ii)(c-d\ii)=(ac-bd)-(ad+bc)\ii .
  \]

  (b) $z\bar z=(a+b\ii)(a-b\ii)=a^2-b^2\ii^2=a^2+b^2=\abs{z}^2$.

  (c) By (b) and (a), $\abs{zw}^2=(zw)\,\overline{zw}=(z\bar z)(w\bar w)=\abs{z}^2\abs{w}^2$.
  Both sides of the claim are non-negative, so we may take square roots.

  (d) $\abs{z}=0$ means $a^2+b^2=0$, which forces $a=b=0$.

  (e) Divide the identity in (b) by $\abs{z}^2\neq0$: $z\cdot(\bar z/\abs{z}^2)=1$.
\end{proof}

Part (e) is how inverses are computed in practice: \emph{multiply top and
bottom by the conjugate.}

\begin{example}\label{ex:inverse}
  Let $z=2-3\ii$. Compute $\bar z$, $\abs{z}$, $z^{-1}$ and $z^2$, and
  simplify $\dfrac{2+\ii}{1-3\ii}$.
\end{example}
\begin{solution}
  Here $\bar z=2+3\ii$ and $\abs{z}^2=4+9=13$, so $\abs{z}=\sqrt{13}$. By
  Theorem~\ref{thm:basic}(e),
  \[
    z^{-1}=\frac{2+3\ii}{13}.
  \]
  Next, $z^2=4-12\ii+9\ii^2=-5-12\ii$. \emph{Check:} by part (c),
  $\abs{z^2}$ must be $\abs{z}^2=13$, and indeed $\sqrt{25+144}=13$.

  For the quotient, multiply numerator and denominator by the conjugate of the
  denominator:
  \[
    \frac{2+\ii}{1-3\ii}
    =\frac{(2+\ii)(1+3\ii)}{(1-3\ii)(1+3\ii)}
    =\frac{2+6\ii+\ii+3\ii^2}{1+9}
    =\frac{-1+7\ii}{10}.
  \]
\end{solution}

\begin{pitfall}
  $\abs{z}^2$ and $z^2$ are different things. For $z=\ii$ we get $\abs{\ii}^2=1$
  but $\ii^2=-1$. The correct statement is $\abs{z}^2=z\bar z$.
\end{pitfall}

The modulus also obeys the inequality that makes it a good notion of size.

\begin{theorem}[Triangle inequality]\label{thm:triangle}
  For all $z,w\in\C$, $\ \abs{z+w}\le\abs{z}+\abs{w}$.
\end{theorem}
\begin{proof}
  Expanding with Theorem~\ref{thm:basic}(a) and (b),
  \[
    \abs{z+w}^2=(z+w)(\bar z+\bar w)=\abs{z}^2+\abs{w}^2+2\Re(z\bar w),
  \]
  because $z\bar w+\bar z w=z\bar w+\overline{z\bar w}=2\Re(z\bar w)$. Now
  $\Re(u)\le\abs{u}$ for every $u$, and $\abs{z\bar w}=\abs{z}\abs{w}$ by
  Theorem~\ref{thm:basic}(c), so
  \[
    \abs{z+w}^2\le\abs{z}^2+\abs{w}^2+2\abs{z}\abs{w}=(\abs{z}+\abs{w})^2 .
  \]
  Taking square roots gives the result.
\end{proof}

\begin{remark}
  The proof used one trick worth remembering: to estimate a modulus, estimate
  its \emph{square}, because $\abs{u}^2=u\bar u$ turns it into algebra. The same
  idea will prove the Cauchy--Schwarz inequality for vectors in Section~\ref{sec:hilbert}.
\end{remark}

\subsection{The exponential and the polar form}

Cartesian form is best for adding. For multiplying we want a better
description, and it comes from the exponential function.

\begin{definition}[Exponential]\label{def:exp}
  For $z\in\C$, $\displaystyle e^z=\exp(z)=\sum_{n=0}^{\infty}\frac{z^n}{n!}$.
\end{definition}

The series converges for every $z$. We use without proof the property that
makes the exponential useful: $e^{z+w}=e^ze^w$ for all $z,w\in\C$ (it follows by
multiplying out the two series). The case of an imaginary exponent is the
celebrated formula of Euler.

\begin{theorem}[Euler's formula]\label{thm:euler}
  For every $\theta\in\R$, $\ e^{\ii\theta}=\cos\theta+\ii\sin\theta$.
\end{theorem}
\begin{proof}
  Put $z=\ii\theta$ in the series. The powers of $\ii$ are $\ii^{2k}=(-1)^k$ and
  $\ii^{2k+1}=(-1)^k\ii$. The series converges absolutely, so we may split it
  into even and odd terms:
  \[
    e^{\ii\theta}
    =\sum_{k=0}^{\infty}\frac{(-1)^k\theta^{2k}}{(2k)!}
    +\ii\sum_{k=0}^{\infty}\frac{(-1)^k\theta^{2k+1}}{(2k+1)!}
    =\cos\theta+\ii\sin\theta,
  \]
  by the Taylor series of cosine and sine.
\end{proof}

In particular $\abs{e^{\ii\theta}}=\sqrt{\cos^2\theta+\sin^2\theta}=1$: the point
$e^{\ii\theta}$ lies on the unit circle, at angle $\theta$ (in radians) from the
positive real axis. Dividing a non-zero $z$ by its modulus puts it on the unit circle, which gives:

\begin{definition}[Polar form]\label{def:polar}
  Every $z\ne0$ can be written $z=r\,e^{\ii\theta}$ with $r=\abs{z}>0$ and
  $\theta\in\R$. The angle $\theta$ is an \term{argument} of $z$. It is determined
  only up to adding multiples of $2\pi$; the one in $(-\pi,\pi]$ is the
  \term{principal argument}, written $\operatorname{Arg}z$.
\end{definition}

The exponential law turns multiplication into the addition of angles, and
multiplying moduli (Figure~\ref{fig:polar}). So, multiplying by $e^{\ii\varphi}$ is a rotation
by $\varphi$.

\begin{figure}[htb]
  \centering
  \begin{tikzpicture}[>=stealth]
    \draw[->,ink] (-0.5,0) -- (5.2,0);
    \draw[->,ink] (0,-0.5) -- (0,3.6);
    \draw[ink!30] (0,0) circle (1);
    \draw[very thick,navy,->] (0,0) -- (25:2.1) node[right] {$z=r\,e^{\ii\theta}$};
    \draw[very thick,navy!60,->] (0,0) -- (55:1.5) node[above left] {$w=s\,e^{\ii\varphi}$};
    \draw[very thick,figacc,->] (0,0) -- (80:3.15) node[above right] {$zw=rs\,e^{\ii(\theta+\varphi)}$};
    \draw[navy] (0.9,0) arc (0:25:0.9) node[midway,right,font=\footnotesize] {$\theta$};
    \draw[navy!60] (25:1.25) arc (25:55:1.25) node[midway,right=1pt,font=\footnotesize] {$\varphi$};
  \end{tikzpicture}
  \caption{Multiplication in polar form: the moduli multiply and the angles add.}
  \label{fig:polar}
\end{figure}

\begin{proposition}[Polar arithmetic]\label{prop:polar}
  If $z=r\,e^{\ii\theta}$ and $w=s\,e^{\ii\varphi}$, then
  \[
    zw=rs\,e^{\ii(\theta+\varphi)},\qquad
    \bar z=r\,e^{-\ii\theta},\qquad
    z^{-1}=r^{-1}e^{-\ii\theta},\qquad
    z^n=r^n e^{\ii n\theta}\ (n\in\Z).
  \]
\end{proposition}
\begin{proof}
  The first formula is the exponential law. For the second, $\overline{e^{\ii\theta}}=\cos\theta-\ii\sin\theta=e^{-\ii\theta}$.
  The third follows from the first with $w=r^{-1}e^{-\ii\theta}$, since then $zw=1$. The last is the first, repeated.
\end{proof}

This lets us compute things that would be painful in Cartesian form.

\begin{example}\label{ex:trig}
  Derive the addition formulas for $\sin(\alpha+\beta)$ and $\cos(\alpha+\beta)$
  from the exponential law.
\end{example}
\begin{solution}
  Expand $e^{\ii(\alpha+\beta)}=e^{\ii\alpha}e^{\ii\beta}$ in two ways. On the left,
  Euler's formula gives $\cos(\alpha+\beta)+\ii\sin(\alpha+\beta)$. On the right,
  \[
    (\cos\alpha+\ii\sin\alpha)(\cos\beta+\ii\sin\beta)
    =(\cos\alpha\cos\beta-\sin\alpha\sin\beta)+\ii(\sin\alpha\cos\beta+\cos\alpha\sin\beta).
  \]
  For real $\alpha,\beta$ both sides have real and imaginary parts that must match, so
  \begin{align*}
    \cos(\alpha+\beta)&=\cos\alpha\cos\beta-\sin\alpha\sin\beta,\\
    \sin(\alpha+\beta)&=\sin\alpha\cos\beta+\cos\alpha\sin\beta.
  \end{align*}
\end{solution}

\begin{example}\label{ex:power}
  Compute $(1+\ii)^{20}$.
\end{example}
\begin{solution}
  Expanding is hopeless, so use the polar form. We have $\abs{1+\ii}=\sqrt2$ and
  the angle is $\pi/4$, so $1+\ii=\sqrt2\,e^{\ii\pi/4}$. By Proposition~\ref{prop:polar},
  \[
    (1+\ii)^{20}=(\sqrt2)^{20}e^{\ii\,5\pi}=2^{10}e^{\ii\pi}=-1024 .
  \]
  \emph{Check:} $(1+\ii)^4=-4$ and $(-4)^5=-1024$.
\end{solution}

\subsubsection*{Finding the argument}

To go from $a+b\ii$ to polar form we need $\theta$. The formula $\theta=\tan^{-1}(b/a)$
is tempting, but $\tan^{-1}$ only returns angles in $(-\pi/2,\pi/2)$, which are
the points with $a>0$. In the other quadrants it is wrong.

\begin{pitfall}
  For $z=-1+\ii$ we have $b/a=-1$ and $\tan^{-1}(-1)=-\pi/4$. But $-1+\ii$ lies in
  the second quadrant, so its angle is $3\pi/4$. The naive formula
  gives the argument of $1-\ii$ instead. \emph{Always look at the signs of $a$ and $b$.}
\end{pitfall}

The correct rule, for $z=a+b\ii\ne0$, is
\[
  \operatorname{Arg}z=
  \begin{cases}
    \tan^{-1}(b/a) & a>0,\\
    \tan^{-1}(b/a)+\pi & a<0,\ b\ge0,\\
    \tan^{-1}(b/a)-\pi & a<0,\ b<0,\\
    \pi/2 & a=0,\ b>0,\\
    -\pi/2 & a=0,\ b<0.
  \end{cases}
\]
This is the function \texttt{atan2(b, a)} found in every programming language.

\begin{example}\label{ex:args}
  Find the polar form of $1+\ii$, $-1+\ii$, $-1-\ii$ and $1-\ii$.
\end{example}
\begin{solution}
  All four have modulus $\sqrt2$, and $\tan^{-1}(b/a)=\pm\pi/4$ in each case; we correct it by the rule above.
  \begin{center}
  \begin{tabular}{cccc}
    \toprule
    $z$ & signs of $(a,b)$ & $\operatorname{Arg}z$ & polar form\\
    \midrule
    $1+\ii$  & $(+,+)$ & $\ \ \pi/4$    & $\sqrt2\,e^{\ii\pi/4}$\\
    $-1+\ii$ & $(-,+)$ & $-\pi/4+\pi=3\pi/4$   & $\sqrt2\,e^{3\ii\pi/4}$\\
    $-1-\ii$ & $(-,-)$ & $\ \ \pi/4-\pi=-3\pi/4$ & $\sqrt2\,e^{-3\ii\pi/4}$\\
    $1-\ii$  & $(+,-)$ & $-\pi/4$       & $\sqrt2\,e^{-\ii\pi/4}$\\
    \bottomrule
  \end{tabular}
  \end{center}
  The four points are the corners of a square, which is a quick sanity check.
\end{solution}

A longer example, where Cartesian form is awkward and polar form is not:

\begin{example}\label{ex:z10}
  Let $z=1+(\sqrt2-1)\ii$. Write $z$ in polar form and compute $w=z^{10}$.
  Which problem appears when recovering the argument of $w$?
\end{example}
\begin{solution}
  We use the known value $\tan(\pi/8)=\sqrt2-1$. Then
  \[
    z=1+\tan\tfrac{\pi}{8}\,\ii
    =\frac{1}{\cos\frac{\pi}{8}}\Bigl(\cos\tfrac{\pi}{8}+\ii\sin\tfrac{\pi}{8}\Bigr)
    =\frac{1}{\cos\frac{\pi}{8}}\,e^{\ii\pi/8}.
  \]
  So $\abs{z}=1/\cos\frac{\pi}{8}$ and $\operatorname{Arg}z=\pi/8$.
  (\emph{Check:} $\cos^2\frac\pi8=\frac{1+\cos(\pi/4)}{2}=\frac{2+\sqrt2}{4}$, hence
  $\abs{z}^2=\frac{4}{2+\sqrt2}=4-2\sqrt2$, and directly
  $\abs{z}^2=1+(\sqrt2-1)^2=4-2\sqrt2$.)

  By Proposition~\ref{prop:polar},
  \[
    w=z^{10}=\abs{z}^{10}e^{\ii\,10\pi/8}=(4-2\sqrt2)^5\,e^{\ii\,5\pi/4}\approx 2.207\,e^{\ii\,5\pi/4}.
  \]
  The problem: $5\pi/4$ is not in $(-\pi,\pi]$. The principal argument is
  $5\pi/4-2\pi=-3\pi/4$, since $w$ lies in the third quadrant. And
  $\tan^{-1}(\Im w/\Re w)=\tan^{-1}(1)=\pi/4$ is wrong, exactly as in the pitfall above.
  In Cartesian form, $w=-\tfrac{\sqrt2}{2}(4-2\sqrt2)^5(1+\ii)\approx-1.561(1+\ii)$.
\end{solution}

\subsection{Phases}

The numbers of modulus one deserve a name.

\begin{definition}[Phase]\label{def:phase}
  A \term{phase} is a complex number of modulus $1$, that is, $e^{\ii\theta}$ for some $\theta\in\R$.
\end{definition}

By Proposition~\ref{prop:polar}, a product of phases is a phase, and so is the
inverse of a phase. Multiplying by a phase \emph{rotates} a number without
changing its modulus: $\abs{e^{\ii\theta}z}=\abs{z}$. The question
that will matter in quantum mechanics is what happens when we rotate
\emph{different} numbers by \emph{different} angles. This is exactly what the
double slit does.

\begin{example}\label{ex:relative-phase}
  Let $z_1=z_2=1/\sqrt2$. Compute $\abs{e^{\ii\theta_1}z_1+e^{\ii\theta_2}z_2}^2$,
  and show that it depends on $\theta_1,\theta_2$ only through $\theta_1-\theta_2$.
\end{example}
\begin{solution}
  Put $u=e^{\ii\theta_1}z_1+e^{\ii\theta_2}z_2$. Then
  $u\bar u=\tfrac12\bigl(e^{\ii\theta_1}+e^{\ii\theta_2}\bigr)\bigl(e^{-\ii\theta_1}+e^{-\ii\theta_2}\bigr)$, which expands to
  \[
    \abs{u}^2=\tfrac12\bigl(1+1+e^{\ii(\theta_1-\theta_2)}+e^{-\ii(\theta_1-\theta_2)}\bigr)
    =1+\cos(\theta_1-\theta_2),
  \]
  using $e^{\ii x}+e^{-\ii x}=2\cos x$. This depends only on the difference $\theta_1-\theta_2$.
  If both numbers are rotated by the same angle ($\theta_1=\theta_2$) it equals $2$ for every angle: that
  common rotation is invisible. If $\theta_1-\theta_2=\pi$ it equals $0$: the two terms cancel.
\end{solution}

\begin{xref}
  In Section~\ref{sec:measurement} the quantities $z_i$ will be amplitudes of a
  qubit. A common phase on \emph{all} of them is the \emph{global phase},
  which cannot be observed. A phase on \emph{some} of them is a \emph{relative
  phase}, and it changes measurement outcomes, as Example~\ref{ex:relative-phase} shows.
\end{xref}

%----------------------------------------------------------------------
\begin{exerciseset}
  \exercise{1} Write in the form $a+b\ii$: $(1+\ii)^2$, $\ \ii(2-3\ii)^2$, $\ (1-\ii)^3$, $\ (3-2\ii)+(1+5\ii)$.

  \exercise{1} Compute $\ii^{7}$, $\ii^{42}$, $\ii^{-3}$ and $\ii^{2026}$. Prove that
  $\ii^n=\ii^{\,n \bmod 4}$ for every $n\in\Z$, including negative $n$.

  \exercise{1} Solve in $\C$: $z^2=-9$, $\ z^2+2z+5=0$, $\ z^2-(1+\ii)z+\ii=0$. Which of these have no real roots?

  \exercise{1} Compute $(3+4\ii)^{-1}$, $(1-\ii)^{-1}$ and $\dfrac{1}{\ii}-\dfrac{3-\ii}{2+\ii}$.

  \exercise{2} Prove that $z\in\R\iff z=\bar z$, that $z\in\ii\R\iff z=-\bar z$, and that
  $\Re z=\frac{z+\bar z}{2}$, $\ \Im z=\frac{z-\bar z}{2\ii}$.

  \exercise{2} Prove that the multiplicative inverse of a non-zero $z\in\C$ is unique.
  Deduce $\abs{z^{-1}}=\abs{z}^{-1}$ and $\abs{z^n}=\abs{z}^n$ for all $n\in\Z$.

  \exercise{2} Prove the reverse triangle inequality $\bigl|\abs{z}-\abs{w}\bigr|\le\abs{z-w}$.
  When does equality $\abs{z+w}=\abs{z}+\abs{w}$ hold in the triangle inequality?

  \exercise{2} Describe and sketch $\{\abs{z-\ii}<2\}$, $\ \{\abs{z-1}=\abs{z+1}\}$ and $\{\Re(z^2)=0\}$.

  \exercise{2} Interpret geometrically $z\mapsto\bar z$, $\ z\mapsto\ii z$, $\ z\mapsto z+w_0$ and $z\mapsto 2z$ as maps of the plane.
  Which are linear over $\R$ (that is, linear maps $\R^2\to\R^2$)? Which satisfy $f(\lambda z)=\lambda f(z)$ for all complex $\lambda$?

  \exercise{1} Compute in Cartesian form: $e^{\ii\pi/2}$, $\ e^{-\ii\pi/3}$, $\ e^{\ii 3\pi/2}$, $\ e^{\ln 2+\ii\pi/6}$.

  \exercise{1} Find modulus and principal argument of $-\sqrt3+\ii$, $\ -2$, $\ -3\ii$ and $\tfrac12-\tfrac{\sqrt3}{2}\ii$.

  \exercise{2} Compute $(\sqrt3+\ii)^9$ and $\left(\dfrac{1+\ii}{1-\ii}\right)^{2027}$.

  \exercise{2} Solve $z^3=-8$, $\ z^2=\ii$ and $\ z^6=-1$. Sketch the solutions: what shape do they form?

  \exercise{2} Use $e^{\ii3\theta}=(e^{\ii\theta})^3$ to express $\cos3\theta$ and $\sin3\theta$ in terms of $\cos\theta$ and $\sin\theta$.

  \exercise{2} Let $U(1)=\{z\in\C:\abs{z}=1\}$. Show that $U(1)$ is closed under products and inverses.
  Show that $\abs{e^{\ii\varphi}z}=\abs{z}$ for all $z$, and that $\abs{e^{\ii\varphi_1}z_1+e^{\ii\varphi_2}z_2}$
  \emph{does} depend on $\varphi_1-\varphi_2$ for general $z_1,z_2$.

  \exercise{3} Let $\alpha,\beta\in\C$ with $\abs{\alpha}^2+\abs{\beta}^2=1$. Show that
  $\abs{\alpha},\abs{\beta}\le1$, that $\abs{\alpha+\beta}^2\le2$, and that $\abs{\alpha\beta}\le\tfrac12$.
  When are the last two bounds attained?

  \exercise{3} Show that $e^z=-1$ exactly when $z=\ii\pi(2k+1)$ with $k\in\Z$. Then find all values of
  $\ii^{\ii}$, defined as $e^{\ii\log\ii}$ using every logarithm of $\ii$. Are they real?
\end{exerciseset}

%======================================================================
\section{Vector spaces and linear maps}\label{sec:vector-spaces}
%======================================================================

\begin{objectives}
  You will decide whether a set is a vector space, find bases and dimensions, recognise linear maps and compute their kernel and range, write a linear map as a matrix and change its basis, and use the Pauli matrices, the first quantum operators we meet.
\end{objectives}

The Introduction stressed that Schr\"odinger's equation is linear: its solutions can be added and multiplied by complex numbers, and that is why amplitudes add. The natural setting for such objects is a vector space, and the maps that respect addition and scaling are the linear maps. A great deal of quantum mechanics is just the study of these two notions, so we spend some time on them.

\subsection{Vector spaces}

Throughout, the scalars $F$ will be either $\R$ or $\C$; in quantum mechanics it is nearly always $\C$.

\begin{definition}[Vector space]\label{def:vs}
  A \term{vector space} over $F$ is a set $V$ with an addition $V\times V\to V$ and a scalar multiplication $F\times V\to V$ such that $(V,+)$ is a commutative group, with neutral element $0_V$, and for all $\lambda,\mu\in F$ and $v,w\in V$:
  \[
    \lambda(\mu v)=(\lambda\mu)v,\quad (\lambda+\mu)v=\lambda v+\mu v,\quad \lambda(v+w)=\lambda v+\lambda w,\quad 1\,v=v .
  \]
\end{definition}

Elements of $V$ are called \term{vectors}. In words: you may add vectors, scale them, and the usual rules of arithmetic hold. From these axioms one deduces that $0\,v=0_V$ for every $v$.

The most important example is $F^n$, the columns of $n$ scalars. But the definition covers much more. The $m\times n$ matrices form a vector space, and so do the polynomials of degree at most $n$, written $P_n$. In infinite dimension we have the continuous functions on $[0,1]$, and, because the equation is linear, the solutions of the Schr\"odinger equation. Sequences $(z_1,z_2,\dots)$ of complex numbers with termwise operations form a vector space as well.

To see how the definition is used, and how it fails, consider two subsets of $\R^2$.

\begin{example}\label{ex:vs-nonexample}
  Are $A=\{(x,y)\in\R^2: x+y=0\}$ and $B=\{(x,y)\in\R^2:x+y=1\}$ vector spaces with the usual operations?
\end{example}
\begin{solution}
  \begin{steps}
    \item \emph{$A$ contains $0$.} $(0,0)\in A$.
    \item \emph{$A$ is closed under sums.} If $x+y=0$ and $x'+y'=0$ then $(x+x')+(y+y')=0$.
    \item \emph{$A$ is closed under scalars.} $\lambda x+\lambda y=\lambda(x+y)=0$. The axioms are inherited from $\R^2$, so $A$ is a vector space (a line through the origin).
    \item \emph{$B$ has no zero.} $(0,0)\notin B$, so $B$ cannot be a vector space.
    \item \emph{$B$ fails closure.} $(1,0)$ and $(0,1)$ are in $B$, but their sum $(1,1)$ is not.
  \end{steps}
\end{solution}

The test used for $A$ is the one to remember: to check that a subset of a known vector space is itself a vector space, it is enough to check that it contains $0$ and is closed under sums and scalar multiples. Such a subset is called a \term{subspace}.

\begin{example}[Is the set of $2\times2$ trace-zero matrices a subspace?]\label{ex:trace-zero}
  The \term{trace} of a square matrix is the sum of its diagonal entries, $\Tr M=\sum_iM_{ii}$. Decide whether $W=\{M\in M_2(\C):\Tr M=0\}$ is a subspace of $M_2(\C)$, and find its dimension.
\end{example}
\begin{solution}
  \begin{steps}
    \item \emph{Subspace test.} $\Tr0=0$; and $\Tr(M+N)=\Tr M+\Tr N=0$, $\Tr(\lambda M)=\lambda\Tr M=0$ for $M,N\in W$. So $W$ is a subspace.
    \item \emph{Describe its elements.} $M=\begin{pmatrix}a&b\\c&-a\end{pmatrix}=a\begin{pmatrix}1&0\\0&-1\end{pmatrix}+b\begin{pmatrix}0&1\\0&0\end{pmatrix}+c\begin{pmatrix}0&0\\1&0\end{pmatrix}$.
    \item \emph{Dimension.} The three matrices are independent and generate $W$, so $\dim W=3$ (compare with $\dim M_2(\C)=4$: one linear condition lowers the dimension by one). We will recognise a nicer basis, $\sx,\sy,\sz$, in Theorem~\ref{thm:paulibasis}.
  \end{steps}
\end{solution}

\begin{pitfall}
  The set of qubit states is \emph{not} a vector space. A qubit state is a \emph{unit} vector of $\C^2$, and the sum of two unit vectors is generally not a unit vector, nor is twice a unit vector. The states live \emph{inside} a vector space ($\C^2$), but they do not form one. This is why normalisation will keep reappearing.
\end{pitfall}

\subsection{Span, independence, bases and dimension}

Fix a vector space $V$. A \term{linear combination} of vectors $v_1,\dots,v_k$ is a sum $\sum_i\lambda_iv_i$ with scalars $\lambda_i$, and the set of all of them is the \term{span} of $v_1,\dots,v_k$. Two properties are of interest.

\begin{definition}[Independence, generators, basis]\label{def:basis}
  Vectors $v_1,\dots,v_k$ are \term{linearly independent} if $\sum_i\lambda_iv_i=0$ holds only when all $\lambda_i=0$. They are a \term{system of generators} of $V$ if their span is all of $V$. They are a \term{basis} if they are both. A space with a finite system of generators is \term{finite-dimensional}.
\end{definition}

Independence says that no vector is redundant; generating says that nothing is missing. A basis is the perfect balance, and it is what lets us use coordinates.

\begin{theorem}[Coordinates]\label{thm:coords}
  If $v_1,\dots,v_n$ is a basis of $V$, every $v\in V$ can be written as $v=\sum_i\lambda_iv_i$ in exactly one way. The numbers $\lambda_i$ are the \term{coordinates} of $v$ in this basis.
\end{theorem}
\begin{proof}
  \begin{steps}
    \item \emph{Existence} is the statement that the $v_i$ generate.
    \item \emph{Uniqueness.} Suppose $\sum_i\lambda_iv_i=\sum_i\mu_iv_i$. Subtracting, $\sum_i(\lambda_i-\mu_i)v_i=0$.
    \item \emph{Independence} forces $\lambda_i=\mu_i$ for every $i$. \qedhere
  \end{steps}
\end{proof}

A fundamental fact, which we use without proof, is that in a finite-dimensional space every basis has the same number of elements. That number is the \term{dimension} $\dim V$. Moreover, any $\dim V$ independent vectors, or any $\dim V$ generating vectors, already form a basis, and a subspace $U\subseteq V$ with $\dim U=\dim V$ equals $V$.
For instance $\dim F^n=n$, $\dim P_n=n+1$ (a basis is $1,x,\dots,x^n$) and the $m\times n$ matrices have dimension $mn$.

From now on, in $\C^2$ we write
\[
  \ket0=\begin{pmatrix}1\\0\end{pmatrix},\quad \ket1=\begin{pmatrix}0\\1\end{pmatrix},\quad \ket{\pm}=\frac{\ket0\pm\ket1}{\sqrt2},\quad \ket{\pm\ii}=\frac{\ket0\pm\ii\ket1}{\sqrt2}.
\]
The symbol $\ket{\,\cdot\,}$ is just a way of writing a vector; its meaning will be explained in Section~\ref{sec:hilbert}. The first basis is the \emph{computational basis}.

\begin{example}\label{ex:qubit-bases}
  Show that $\{\ket+,\ket-\}$ is a basis of $\C^2$, write $\ket0$ in it, and write $\ket{+\ii}$ in it.
\end{example}
\begin{solution}
  \begin{steps}
    \item \emph{Independence.} If $\alpha\ket++\beta\ket-=0$, then in components $(\alpha+\beta)/\sqrt2=0$ and $(\alpha-\beta)/\sqrt2=0$, so $\alpha=\beta=0$.
    \item \emph{Basis.} Two independent vectors in a space of dimension $2$ form a basis.
    \item \emph{Express $\ket0,\ket1$.} Adding the definitions, $\ket++\ket-=\sqrt2\,\ket0$; subtracting, $\ket+-\ket-=\sqrt2\,\ket1$. So $\ket0=\tfrac1{\sqrt2}(\ket++\ket-)$ and $\ket1=\tfrac1{\sqrt2}(\ket+-\ket-)$.
    \item \emph{Express $\ket{+\ii}$.} $\ket{+\ii}=\tfrac1{\sqrt2}\bigl(\ket0+\ii\ket1\bigr)=\tfrac12\bigl((1+\ii)\ket++(1-\ii)\ket-\bigr)$.
    \item \emph{Check.} $\tfrac14\bigl(\abs{1+\ii}^2+\abs{1-\ii}^2\bigr)=1$, as it should be for a unit vector (we will see why in Section~\ref{sec:hilbert}).
  \end{steps}
\end{solution}

\begin{example}[A dependent set of functions]\label{ex:dependent-functions}
  Are $\sin x$, $\cos x$ and $\sin(x+1)$ linearly independent as functions on $\R$?
\end{example}
\begin{solution}
  \begin{steps}
    \item \emph{Expand the third function} with the addition formula (Example~\ref{ex:trig}): $\sin(x+1)=\cos1\,\sin x+\sin1\,\cos x$.
    \item \emph{Rearrange.} $\cos1\cdot\sin x+\sin1\cdot\cos x-1\cdot\sin(x+1)=0$: a vanishing combination with coefficients not all zero.
    \item \emph{Conclude.} The three functions are linearly \emph{dependent}. In fact $\Span\{\sin x,\cos x,\sin(x+1)\}$ is two-dimensional.
  \end{steps}
\end{solution}

\begin{example}[Why quantum computers are hard to simulate]\label{ex:nqubits}
  The state space of $n$ qubits is $(\C^2)^{\otimes n}=\C^{2^n}$ (a tensor product, introduced in Section~\ref{sec:tensor}; for now we only need its dimension, $2^n$). How much memory does a general state of $50$ qubits need, at 16 bytes per complex number?
\end{example}
\begin{solution}
  \begin{steps}
    \item \emph{Count the amplitudes.} A general state is a list of $2^{50}\approx1.1\times10^{15}$ complex numbers.
    \item \emph{Memory.} $2^{50}\cdot16=2^{54}\approx1.8\times10^{16}$ bytes, about $18$ petabytes.
    \item \emph{Scale up.} For $300$ qubits the number $2^{300}\approx2\times10^{90}$ exceeds the number of atoms in the observable universe.
    \item \emph{Moral.} The dimension grows exponentially with the number of qubits, and this is the source both of the power of quantum computing and of the difficulty of simulating it.
  \end{steps}
\end{solution}

\subsection{Sums of subspaces}

Two subspaces $W_1,W_2$ of $V$ can be combined: their \term{sum} $W_1+W_2=\{w_1+w_2\}$ is the smallest subspace containing both. The sum is \term{direct}, written $W_1\oplus W_2$, if every vector of the sum can be written \emph{uniquely} as $w_1+w_2$. This happens exactly when $W_1\cap W_2=\{0\}$. If $V=W_1\oplus W_2$, then $\dim V=\dim W_1+\dim W_2$.

A familiar example: every function $f:\R\to\R$ splits as an even part $\tfrac12(f(x)+f(-x))$ plus an odd part $\tfrac12(f(x)-f(-x))$, and the only function that is both even and odd is $0$, so $\{\text{functions}\}=\{\text{even}\}\oplus\{\text{odd}\}$. We will meet this structure again with projectors, which split a space in the same way.

\subsection{Linear maps}

The maps between vector spaces that matter are those that respect the structure.

\begin{definition}[Linear map]\label{def:linmap}
  Let $V,W$ be vector spaces over $F$. A map $T:V\to W$ is \term{linear} if $T(\lambda v)=\lambda T(v)$ and $T(v+w)=T(v)+T(w)$ for all $\lambda\in F$ and $v,w\in V$. The set of linear maps from $V$ to $W$ is written $L(V,W)$, and $L(V)=L(V,V)$.
\end{definition}

A linear map sends $0$ to $0$, so any map with $T(0)\neq0$ is not linear. The map $T(x,y)=2x+y$ is linear, while $T(x,y)=2x+y-1$ is not, because $T(0,0)=-1$. Differentiation $f\mapsto f'$ and integration $f\mapsto\int_0^1f$ are linear; squaring $f\mapsto f^2$ is not. A linear map is determined by what it does to a basis: if $v=\sum_i\lambda_iv_i$ then $T(v)=\sum_i\lambda_iT(v_i)$.

Linear maps can themselves be added and scaled, so $L(V,W)$ is a vector space, of dimension $\dim V\cdot\dim W$ in finite dimension. Maps can also be composed, and composition of linear maps is linear. It is not commutative in general, and this non-commutativity will be at the heart of quantum mechanics.

\begin{example}[Composition does not commute]\label{ex:noncommuting}
  On $\R^2$ let $T(x,y)=(y,x)$ (swap) and $S(x,y)=(x,0)$ (project to the $x$-axis). Compute $ST$ and $TS$.
\end{example}
\begin{solution}
  \begin{steps}
    \item \emph{$ST$ means: first $T$, then $S$.} $ST(x,y)=S(y,x)=(y,0)$.
    \item \emph{$TS$ means: first $S$, then $T$.} $TS(x,y)=T(x,0)=(0,x)$.
    \item \emph{Compare.} $ST(1,0)=(0,0)$ but $TS(1,0)=(0,1)$: the two compositions differ, and $ST-TS\neq0$.
  \end{steps}
\end{solution}

\begin{lookahead}
  The operator $ST-TS$ is called the \term{commutator} of $S$ and $T$ and written $[S,T]$ (and $\{S,T\}=ST+TS$ is their \term{anticommutator}); operators with $[S,T]=0$ \emph{commute}. The size of the commutator is what the uncertainty principle will measure. Observables are operators; two observables that do not commute (Section~\ref{sec:measurement}) cannot both be sharp in the same state, and the position and momentum operators satisfy the extreme version $\hat x\hat p-\hat p\hat x=\ii\hbar$.
\end{lookahead}

Two subspaces measure what a linear map $T:V\to W$ does.

\begin{definition}[Kernel and range]\label{def:kerran}
  The \term{kernel} (or null space) of $T$ is $\Ker T=\{v\in V:T(v)=0\}$, a subspace of $V$. The \term{range} of $T$ is $\Ran T=\{T(v):v\in V\}$, a subspace of $W$. We say $T$ is \term{injective} if $T(v_1)=T(v_2)$ implies $v_1=v_2$, and \term{surjective} if $\Ran T=W$.
\end{definition}

The kernel is the information that $T$ destroys; the range is what it can reach. The dimension of the range is called the \term{rank} of $T$. The two ideas are linked, as the next two results show.

\begin{theorem}\label{thm:inj-ker}
  A linear map $T$ is injective if and only if $\Ker T=\{0\}$.
\end{theorem}
\begin{proof}
  \begin{steps}
    \item \emph{($\Rightarrow$)} If $T$ is injective and $T(v)=0=T(0)$, then $v=0$.
    \item \emph{($\Leftarrow$)} Suppose $\Ker T=\{0\}$ and $T(v_1)=T(v_2)$. Then $T(v_1-v_2)=0$, so $v_1-v_2=0$. \qedhere
  \end{steps}
\end{proof}

\begin{theorem}[Rank--nullity]\label{thm:ranknull}
  If $V$ is finite-dimensional and $T\in L(V,W)$, then $\dim V=\dim\Ker T+\dim\Ran T$.
\end{theorem}
\begin{proof}
  \begin{steps}
    \item \emph{Build a basis adapted to $T$.} Take a basis $u_1,\dots,u_k$ of $\Ker T$ and extend it to a basis $u_1,\dots,u_k,v_1,\dots,v_r$ of $V$. We claim $T(v_1),\dots,T(v_r)$ is a basis of $\Ran T$; this gives $\dim V=k+r$.
    \item \emph{They generate.} For any $v=\sum a_iu_i+\sum b_jv_j$ we have $T(v)=\sum b_jT(v_j)$, since $T(u_i)=0$.
    \item \emph{They are independent.} If $\sum b_jT(v_j)=0$, then $\sum b_jv_j\in\Ker T$, so $\sum b_jv_j=\sum a_iu_i$ for some $a_i$.
    \item \emph{Conclude.} Independence of the whole basis forces all the $a_i$ and $b_j$ to vanish. \qedhere
  \end{steps}
\end{proof}

\begin{corollary}\label{cor:inj-surj}
  If $\dim V<\infty$ and $T\in L(V)$, then $T$ is injective $\iff$ $T$ is surjective $\iff$ $T$ is invertible.
\end{corollary}
\begin{proof}
  \begin{steps}
    \item By Theorem~\ref{thm:inj-ker}, $T$ is injective iff $\dim\Ker T=0$.
    \item By rank--nullity this holds iff $\dim\Ran T=\dim V$, that is, iff $\Ran T=V$. \qedhere
  \end{steps}
\end{proof}

The next example shows that finite dimension cannot be dropped.

\begin{example}\label{ex:shifts}
  On the space of complex sequences $(z_1,z_2,\dots)$ consider the left shift $\ell(z_1,z_2,z_3,\dots)=(z_2,z_3,\dots)$ and the right shift $\rho(z_1,z_2,z_3,\dots)=(0,z_1,z_2,\dots)$. Are they injective, surjective, bijective? Why does this not contradict Corollary~\ref{cor:inj-surj}?
\end{example}
\begin{solution}
  \begin{steps}
    \item \emph{$\ell$ is not injective.} It forgets the first entry: $\ell(1,0,0,\dots)=0$.
    \item \emph{$\ell$ is surjective.} $(w_1,w_2,\dots)=\ell(0,w_1,w_2,\dots)$.
    \item \emph{$\rho$ is injective.} $\rho(z)=0$ forces all $z_i=0$.
    \item \emph{$\rho$ is not surjective.} Nothing maps to $(1,0,0,\dots)$. Neither shift is bijective.
    \item \emph{No contradiction.} The space is infinite-dimensional, and Corollary~\ref{cor:inj-surj} needs rank--nullity, which needs $\dim V<\infty$.
  \end{steps}
\end{solution}

\begin{example}\label{ex:ranknull}
  Let $T:\R^3\to\R^2$, $T(x,y,z)=(x+y,\,y+z)$. Find $\Ker T$, $\Ran T$ and check rank--nullity.
\end{example}
\begin{solution}
  \begin{steps}
    \item \emph{Kernel.} $T(x,y,z)=0$ means $x=-y$ and $z=-y$, so $\Ker T=\{(-t,t,-t):t\in\R\}=\Span\{(-1,1,-1)\}$, of dimension $1$.
    \item \emph{Range.} $T(1,0,0)=(1,0)$ and $T(0,0,1)=(0,1)$, so $\Ran T=\R^2$, of dimension $2$.
    \item \emph{Check.} $\dim\R^3=3=1+2$.
  \end{steps}
\end{solution}

\begin{figure}[htb]
  \centering
  \begin{tikzpicture}[>=stealth]
    \node[font=\small] at (-4.5,1.15) {$V$};
    \draw[fill=ink!8,draw=ink] (-5.7,-0.6) rectangle (-1.3,1.0);
    \draw[fill=figacc!25,draw=figacc] (-5.5,-0.4) rectangle (-3.6,0.8);
    \draw[fill=navy!20,draw=navy] (-3.6,-0.4) rectangle (-1.5,0.8);
    \node[font=\small,figacc] at (-4.55,0.2) {$\Ker T$};
    \node[font=\small,navy] at (-2.55,0.2) {the rest};
    \draw[->,thick,navy] (-2.55,-0.65) -- (0.9,-0.65) node[midway,below,font=\small] {$T$ is injective here};
    \draw[->,thick,figacc] (-4.55,-0.65)  -- ++(0,-0.3);
    \node[font=\small,figacc] at (-4.55,-1.25) {sent to $0$};
    \node[font=\small] at (3.3,1.15) {$W$};
    \draw[fill=ink!8,draw=ink] (1.3,-0.6) rectangle (5.3,1.0);
    \draw[fill=navy!20,draw=navy] (1.5,-0.4) rectangle (3.4,0.8);
    \node[font=\small,navy] at (2.45,0.2) {$\Ran T$};
    \node[font=\footnotesize] at (4.3,0.2) {unreached};
  \end{tikzpicture}
  \caption{Rank--nullity: the part of $V$ that $T$ kills ($\Ker T$) and the part that it sends one-to-one onto $\Ran T$ add up to all of $V$, so $\dim V=\dim\Ker T+\dim\Ran T$.}
  \label{fig:ranknull}
\end{figure}

\subsection{Matrices}

Linear maps between finite-dimensional spaces become matrices once we pick bases. Let $v_1,\dots,v_n$ be a basis of $V$ and $w_1,\dots,w_m$ a basis of $W$. For each $i$ write $T(v_i)=\sum_jM_{ji}w_j$. The \term{matrix of $T$} in these bases is $M=(M_{ji})$, an $m\times n$ array whose $i$th column holds the coordinates of $T(v_i)$:
\[
  \text{coordinates of }T(v)\ =\ M\cdot\text{coordinates of }v .
\]
This is the one idea to hold on to: \emph{the columns of the matrix are the images of the basis vectors.}

\begin{example}\label{ex:matrices}
  Find the matrix of $D:P_3\to P_3$, $D(p)=p'$, in the basis $1,x,x^2,x^3$, and of the rotation $R_\theta$ of $\R^2$ by an angle $\theta$ in the canonical basis.
\end{example}
\begin{solution}
  \begin{steps}
    \item \emph{Images of the basis of $P_3$.} $D(1)=0$, $D(x)=1$, $D(x^2)=2x$, $D(x^3)=3x^2$.
    \item \emph{Write them as columns of coordinates:}
      \[
        M_D=\begin{pmatrix}0&1&0&0\\0&0&2&0\\0&0&0&3\\0&0&0&0\end{pmatrix}.
      \]
    \item \emph{Check.} Taking four derivatives of a cubic gives $0$, and indeed $M_D^4=0$.
    \item \emph{The rotation.} $R_\theta e_1=(\cos\theta,\sin\theta)$ and $R_\theta e_2=(-\sin\theta,\cos\theta)$, so $M_{R_\theta}=\begin{pmatrix}\cos\theta&-\sin\theta\\ \sin\theta&\cos\theta\end{pmatrix}$.
  \end{steps}
\end{solution}

A matrix depends on the bases chosen, but the map does not. If $P$ is the matrix whose columns are the new basis vectors written in the old basis, the matrix of $T\in L(V)$ changes by $M'=P^{-1}MP$. The next example shows why it is worth choosing the basis well.

\begin{example}\label{ex:changebasis}
  Let $T(x,y)=(x+y,\,x+y)$ on $\R^2$. Find its matrix in the canonical basis and in the basis $b_1=(1,1)$, $b_2=(1,-1)$.
\end{example}
\begin{solution}
  \begin{steps}
    \item \emph{Canonical basis.} $T(e_1)=(1,1)$ and $T(e_2)=(1,1)$, so $M=\begin{pmatrix}1&1\\1&1\end{pmatrix}$.
    \item \emph{New basis.} $T(b_1)=(2,2)=2b_1$ and $T(b_2)=(0,0)=0\,b_2$, so $M'=\begin{pmatrix}2&0\\0&0\end{pmatrix}$.
    \item \emph{Check the formula $M'=P^{-1}MP$} with $P=\begin{pmatrix}1&1\\1&-1\end{pmatrix}$ and $P^{-1}=\tfrac12\begin{pmatrix}1&1\\1&-1\end{pmatrix}$: $MP=\begin{pmatrix}2&0\\2&0\end{pmatrix}$ and $P^{-1}MP=\begin{pmatrix}2&0\\0&0\end{pmatrix}$.
    \item \emph{Reading.} In the good basis the map is simply ``stretch $b_1$ by $2$ and kill $b_2$'' (Figure~\ref{fig:changebasis}).
  \end{steps}
\end{solution}

\begin{figure}[htb]
  \centering
  \begin{tikzpicture}[>=stealth,scale=1.15]
    \draw[ink!20] (-2.4,-2.4) grid (2.4,2.4);
    \draw[->,ink] (-2.4,0)--(2.4,0); \draw[->,ink] (0,-2.4)--(0,2.4);
    \draw[very thick,navy,->] (0,0)--(1,1) node[above left] {$b_1$};
    \draw[very thick,navy,->] (0,0)--(1,-1) node[below right] {$b_2$};
    \draw[very thick,figacc,->] (0,0)--(2,2) node[right] {$T(b_1)=2b_1$};
    \fill[figacc] (0,0) circle (2.2pt);
    \node[figacc,font=\small,anchor=north east] at (-0.1,-0.1) {$T(b_2)=0$};
  \end{tikzpicture}
  \caption{The map $T(x,y)=(x+y,x+y)$ in the basis $b_1,b_2$: it stretches $b_1$ by a factor $2$ and sends $b_2$ to $0$. Its matrix there is $\diag(2,0)$.}
  \label{fig:changebasis}
\end{figure}

\paragraph{Rank, transpose and adjoint.}
The \term{rank} of a matrix is the dimension of the space generated by its columns, which is the same as that generated by its rows, and it equals $\dim\Ran T$. The \term{transpose} $M^{t}$ swaps rows and columns, $(M^t)_{ij}=M_{ji}$. The \term{adjoint} (or conjugate transpose) $M^\dagger$ also conjugates each entry, $(M^\dagger)_{ij}=\overline{M_{ji}}$. For real matrices the two coincide.
All three have the same rank, and $(AB)^t=B^tA^t$, $(AB)^\dagger=B^\dagger A^\dagger$: the order flips.

\paragraph{Square matrices.}
Some square matrices $M$ over $\C$ have names that we will use constantly. The \term{trace} is $\Tr M=\sum_iM_{ii}$, the sum of the diagonal.

\begin{definition}[Special matrices]\label{def:special}
  A square complex matrix $M$ is called
  \begin{itemize}
    \item \term{Hermitian} if $M^\dagger=M$;
    \item \term{unitary} if $M^\dagger M=MM^\dagger=\Id$, that is, $M^{-1}=M^\dagger$;
    \item \term{involutory} if $M^2=\Id$;
    \item \term{idempotent} if $M^2=M$, and a \term{projector} if it is also Hermitian;
    \item \term{nilpotent} if $M^k=0$ for some $k$.
  \end{itemize}
\end{definition}

Hermitian matrices will be the observables of quantum mechanics, and unitary matrices the reversible quantum gates. The determinant, which you already know, tells us when a square matrix can be undone: for $A\in M_n(F)$ the following are all equivalent. $A$ is invertible; $\det A\neq0$; $\rank A=n$; the columns of $A$ are independent; the rows of $A$ are independent; $0$ is not an eigenvalue of $A$ (eigenvalues are the subject of Section~\ref{sec:operators}); the system $Ax=b$ has a unique solution for every $b$.

\begin{example}[Recognising special matrices]\label{ex:special}
  Which of $A=\begin{pmatrix}0&1\\0&0\end{pmatrix}$, $B=\tfrac12\begin{pmatrix}1&1\\1&1\end{pmatrix}$, $C=\begin{pmatrix}1&1\\0&0\end{pmatrix}$ are Hermitian, idempotent, projectors, nilpotent?
\end{example}
\begin{solution}
  \begin{steps}
    \item \emph{$A$.} $A^2=0$: nilpotent. $A^\dagger=\begin{pmatrix}0&0\\1&0\end{pmatrix}\neq A$: not Hermitian, so not a projector.
    \item \emph{$B$.} $B^\dagger=B$ and $B^2=\tfrac14\begin{pmatrix}2&2\\2&2\end{pmatrix}=B$: Hermitian and idempotent, hence a projector (onto $\ket+$).
    \item \emph{$C$.} $C^2=\begin{pmatrix}1&1\\0&0\end{pmatrix}=C$: idempotent. But $C^\dagger\neq C$, so it is not a projector: it projects onto the $x$-axis \emph{obliquely}.
  \end{steps}
\end{solution}

\subsection{The Pauli matrices}

Three $2\times2$ complex matrices are the workhorses of quantum computing:
\[
  \sx=\begin{pmatrix}0&1\\1&0\end{pmatrix},\qquad \sy=\begin{pmatrix}0&-\ii\\\ii&0\end{pmatrix},\qquad \sz=\begin{pmatrix}1&0\\0&-1\end{pmatrix}.
\]
They are called the \term{Pauli matrices}; in a quantum circuit they are the gates $X$, $Y$ and $Z$. We write $\sigma_1,\sigma_2,\sigma_3$ for $\sx,\sy,\sz$.

\begin{proposition}\label{prop:pauli}
  Each Pauli matrix is Hermitian, unitary, involutory and traceless, with determinant $-1$. Moreover, for $j,k\in\{1,2,3\}$,
  \[
    \sigma_j\sigma_k=\delta_{jk}\Id+\ii\sum_{l=1}^{3}\varepsilon_{jkl}\,\sigma_l,
  \]
  where $\varepsilon_{jkl}$ is $1$ for $(j,k,l)$ an even permutation of $(1,2,3)$, $-1$ for an odd one, and $0$ if two indices coincide.
\end{proposition}
\begin{proof}
  \begin{steps}
    \item \emph{Hermitian.} Clear for $\sx,\sz$; for $\sy$, swapping rows with columns gives $\begin{pmatrix}0&\ii\\-\ii&0\end{pmatrix}$ and conjugating returns $\sy$.
    \item \emph{Trace and determinant.} The traces are $0$ and the determinants are $-1$, $-(-\ii)(\ii)=-1$, $-1$.
    \item \emph{Products.} $\sx^2=\Id$ and $\sx\sy=\begin{pmatrix}0&1\\1&0\end{pmatrix}\begin{pmatrix}0&-\ii\\\ii&0\end{pmatrix}=\begin{pmatrix}\ii&0\\0&-\ii\end{pmatrix}=\ii\sz$; the other products are computed in the same way.
    \item \emph{Unitary.} Each $\sigma_j$ is Hermitian with $\sigma_j^2=\Id$, so $\sigma_j^\dagger\sigma_j=\Id$. \qedhere
  \end{steps}
\end{proof}

In words: each Pauli matrix squares to the identity, and distinct ones multiply to $\pm\ii$ times the third, with $\sx\sy=\ii\sz$, $\sy\sz=\ii\sx$, $\sz\sx=\ii\sy$ and the opposite sign in the reverse order. In particular $\sigma_j\sigma_k=-\sigma_k\sigma_j$ for $j\neq k$: Pauli matrices \emph{anticommute}. This gives short proofs of identities that look laborious by hand.

\begin{example}\label{ex:hadamard}
  The \emph{Hadamard gate} is $H=\frac{1}{\sqrt2}\begin{pmatrix}1&1\\1&-1\end{pmatrix}$. Show that $H=(\sx+\sz)/\sqrt2$, that $H^2=\Id$, and that $H\sx H=\sz$.
\end{example}
\begin{solution}
  \begin{steps}
    \item \emph{The identity.} $\sx+\sz=\begin{pmatrix}1&1\\1&-1\end{pmatrix}$, which gives the first claim.
    \item \emph{$H^2$.} Use $\sx^2=\sz^2=\Id$ and $\sx\sz+\sz\sx=0$: $H^2=\tfrac12(\sx+\sz)^2=\tfrac12(\sx^2+\sx\sz+\sz\sx+\sz^2)=\tfrac12(\Id+0+\Id)=\Id$.
    \item \emph{$H\sx H$.} $H\sx H=\tfrac12(\sx+\sz)\sx(\sx+\sz)=\tfrac12(\Id+\sz\sx)(\sx+\sz)$.
    \item \emph{Expand.} $\tfrac12\bigl(\sx+\sz+\sz\sx\sx+\sz\sx\sz\bigr)=\tfrac12\bigl(\sx+\sz+\sz-\sx\bigr)=\sz$, because $\sz\sx\sx=\sz$ and $\sz\sx\sz=-\sx\sz\sz=-\sx$.
  \end{steps}
\end{solution}

The identity $H\sx H=\sz$ says that the Hadamard gate exchanges the roles of $\sx$ and $\sz$. We will see why when we study eigenvectors.

\begin{example}[Products of Pauli matrices]\label{ex:pauli-products}
  Compute $\sx\sy\sz$ and $\sy\sx$, and show that $\sx\sy\sz=\ii\Id$.
\end{example}
\begin{solution}
  \begin{steps}
    \item \emph{First product.} $\sx\sy=\ii\sz$ (Proposition~\ref{prop:pauli}).
    \item \emph{Multiply by $\sz$.} $\sx\sy\sz=\ii\sz\sz=\ii\Id$.
    \item \emph{The reversed product.} $\sy\sx=-\sx\sy=-\ii\sz$ by anticommutation.
  \end{steps}
\end{solution}

Finally, the Pauli matrices together with the identity form a particularly useful basis.

\begin{theorem}[Pauli basis]\label{thm:paulibasis}
  $\{\Id,\sx,\sy,\sz\}$ is a basis of $M_2(\C)$ over $\C$. If we write $\sigma_0=\Id$, then every $A\in M_2(\C)$ satisfies $A=\tfrac12\sum_{\mu=0}^{3}\Tr(\sigma_\mu A)\,\sigma_\mu$.
\end{theorem}
\begin{proof}
  \begin{steps}
    \item \emph{Generate.} Let $A=\begin{pmatrix}a&b\\c&d\end{pmatrix}$. A direct check shows $A=\tfrac{a+d}{2}\Id+\tfrac{b+c}{2}\sx+\tfrac{\ii(b-c)}{2}\sy+\tfrac{a-d}{2}\sz$, so the four matrices generate $M_2(\C)$.
    \item \emph{Basis.} Since $\dim M_2(\C)=4$, four generators form a basis.
    \item \emph{Coefficients.} If $A=\sum_\nu c_\nu\sigma_\nu$, multiplying by $\sigma_\mu$ and taking the trace gives $\Tr(\sigma_\mu A)=\sum_\nu c_\nu\Tr(\sigma_\mu\sigma_\nu)=2c_\mu$, because $\Tr(\sigma_\mu\sigma_\nu)=2\delta_{\mu\nu}$ by Proposition~\ref{prop:pauli}, as $\Tr\Id=2$ and the $\sigma_l$ are traceless. \qedhere
  \end{steps}
\end{proof}

\begin{example}\label{ex:paulidecomp}
  Decompose $A=\begin{pmatrix}1&2\\3&-1\end{pmatrix}$ and $\ket0\!\bra1=\begin{pmatrix}0&1\\0&0\end{pmatrix}$ in the Pauli basis.
\end{example}
\begin{solution}
  \begin{steps}
    \item \emph{$A$: read off $a=1,b=2,c=3,d=-1$.} The coefficients are $\tfrac{a+d}{2}=0$, $\tfrac{b+c}{2}=\tfrac52$, $\tfrac{\ii(b-c)}{2}=-\tfrac\ii2$, $\tfrac{a-d}{2}=1$.
    \item \emph{Assemble.} $A=\tfrac52\sx-\tfrac\ii2\sy+\sz$.
    \item \emph{$\ket0\!\bra1$: $a=d=c=0$, $b=1$.} The coefficients give $\tfrac12\sx+\tfrac\ii2\sy=\tfrac12(\sx+\ii\sy)$.
    \item \emph{Check.} $\sx+\ii\sy=\begin{pmatrix}0&2\\0&0\end{pmatrix}$.
  \end{steps}
\end{solution}

\begin{example}[The real dimension of the Hermitian matrices]\label{ex:hermitian-dim}
  Show that the Hermitian $2\times2$ matrices form a real vector space of dimension $4$, with basis $\Id,\sx,\sy,\sz$.
\end{example}
\begin{solution}
  \begin{steps}
    \item \emph{It is a real vector space} (not a complex one): sums and \emph{real} multiples of Hermitian matrices are Hermitian, but $\ii M$ is not Hermitian when $M\neq0$.
    \item \emph{General form.} Hermitian means $a,d\in\R$ and $c=\bar b$, so $M=\begin{pmatrix}a&b\\\bar b&d\end{pmatrix}$, with $b=\beta_1-\ii\beta_2$.
    \item \emph{Pauli expansion.} By Theorem~\ref{thm:paulibasis}, $M=\tfrac{a+d}2\Id+\beta_1\sx+\beta_2\sy+\tfrac{a-d}2\sz$, with all four coefficients real.
    \item \emph{Dimension.} The four matrices are independent over $\R$, so the real dimension is $4$: a Hermitian qubit observable is described by four real numbers, the coefficient of $\Id$ and a vector $\vec\beta\in\R^3$ of coefficients of $\sx,\sy,\sz$. (This is the reason that $\rho=\tfrac12(\Id+\vec r\cdot\vec\sigma)$ will describe all qubit states, Section~\ref{sec:density}.)
  \end{steps}
\end{solution}

\begin{exerciseset}
  \exercise{1} Decide which of these are vector spaces (over $\R$ or $\C$, as indicated): $\{(x,y)\in\R^2:xy\ge0\}$; the $2\times2$ complex matrices of trace $0$; the Hermitian $2\times2$ matrices over $\C$, and over $\R$; the $2\times2$ unitary matrices; polynomials of degree exactly $3$; $\{f:[0,1]\to\R\text{ differentiable}: f(0)=0\}$ and $\{f:f(0)=1\}$.

  \exercise{2} Is the set of unit vectors of $\C^2$ closed under addition? Under multiplication by complex scalars? Under multiplication by a scalar of modulus $1$? Explain what this says about the structure of the set of qubit states.

  \exercise{1} Decide whether these sets are linearly independent: $\{(1,-1),(1,2),(2,1)\}\subset\R^2$; $\{(1,\ii),(\ii,-1)\}\subset\C^2$; $\{1,x,x^2,(1-x)^2\}\subset P_2$; $\{e^x,e^{2x},e^{3x}\}$ as functions on $\R$.

  \exercise{1} Show that $\{\ket0,\ket1\}$, $\{\ket+,\ket-\}$ and $\{\ket{+\ii},\ket{-\ii}\}$ are bases of $\C^2$. Express $\ket0$ in each of them.

  \exercise{2} Prove: any set containing $0$ is dependent; $\{v_1,\dots,v_k\}$ is dependent if and only if some $v_j$ is a linear combination of the others; a set $A$ is a basis if and only if every vector is a \emph{unique} linear combination of elements of $A$.

  \exercise{2} Show that $\{v_1,v_2\}$ is independent if and only if $\{v_1+v_2,\,v_1-v_2\}$ is independent.

  \exercise{2} Let $W_1,W_2$ be subspaces of $V$. Show that $W_1\cap W_2$ is a subspace, give an example in $\R^2$ where $W_1\cup W_2$ is not, and prove that $W_1\cup W_2$ is a subspace only if one is contained in the other.

  \exercise{2} Show that $V=W_1\oplus W_2$ if and only if $V=W_1+W_2$ and $W_1\cap W_2=\{0\}$. Prove that the even and the odd functions are subspaces of $\{f:\R\to\R\}$ and that the space is their direct sum.

  \exercise{1} Compute the dimension of: $\C^n$ over $\C$; $M_{m\times n}(\C)$; $P_n$; the diagonal $n\times n$ matrices; the upper triangular ones; the Hermitian $n\times n$ matrices over $\R$; the traceless $n\times n$ complex matrices.

  \exercise{2} Show that the solutions of $y''+y=0$ form a space of dimension $2$ with basis $\{\cos x,\sin x\}$ (you may assume the existence and uniqueness theorem for ODEs), and find another basis of complex exponentials.

  \exercise{1} Decide which of these maps are linear, and specify domain and codomain: $T(x,y)=(xy,x)$; $T(z)=\bar z$ on $\C$ (over $\R$? over $\C$?); $f\mapsto f+1$; $f\mapsto f(0)$; $A\mapsto A^t$; $A\mapsto\Tr A$; $A\mapsto\det A$; $A\mapsto BA-AB$ for a fixed $B$.

  \exercise{2} Let $\{v_1,\dots,v_n\}$ be a basis of $V$ and $w_1,\dots,w_n\in W$ arbitrary. Prove there is a unique linear $T$ with $T(v_i)=w_i$. Find the linear $T:\R^2\to\R^3$ with $T(1,1)=(1,0,2)$, $T(1,-1)=(0,1,0)$ and compute $T(3,1)$.

  \exercise{2} Prove that $L(V,W)$ is a vector space, build an explicit basis when $\dim V=n$, $\dim W=m$, and conclude that $\dim L(V,W)=mn$.

  \exercise{2} Show that the composition of linear maps is linear and associative. On $\R^2$ let $R$ be the rotation by $90^\circ$ and $F(x,y)=(x,-y)$ the reflection in the $x$-axis; compute $RF$ and $FR$, and describe each geometrically.

  \exercise{2} Let $D:P_4\to P_3$, $D(p)=p'$ and $M:P_3\to P_4$, $M(p)=xp$. Compute $DM-MD$ on $P_3$. Which famous commutation relation does this resemble?

  \exercise{1} For each map find kernel, range and dimensions, and verify rank--nullity: $T(x,y,z)=(x-y,\,y-z)$; $D:P_n\to P_n$, $D(p)=p'$; $\Tr:M_n(\C)\to\C$; $T(A)=A-A^t$ on $M_n(\R)$.

  \exercise{2} Let $T\in L(V)$ with $T^2=0$. Prove $\Ran T\subseteq\Ker T$ and deduce $\dim\Ran T\le\frac12\dim V$. Give an example with equality in $\C^2$.

  \exercise{3} Let $T\in L(V)$ with $T^2=T$. Prove $V=\Ker T\oplus\Ran T$. Find an operator for which $\Ker T\cap\Ran T\neq\{0\}$.

  \exercise{1} Find the matrices of: $T(x,y)=(x+2y,\,3x-y,\,y)$ in the canonical bases; the reflection across $y=x$; the orthogonal projection onto the $x$-axis; $T(A)=A^t$ on $M_2(\C)$ in the basis $E_{11},E_{12},E_{21},E_{22}$.

  \exercise{2} Let $T(x,y)=(2x+y,\,x+2y)$. Find its matrix $M$ in the canonical basis and $M'$ in the basis $b_1=(1,1)$, $b_2=(1,-1)$, and check that $M'=P^{-1}MP$. Then prove in general that $\Tr$, $\det$ and rank are invariant under $M\mapsto P^{-1}MP$.

  \exercise{2} The Hadamard gate is the linear map with $H\ket0=\ket+$ and $H\ket1=\ket-$. Using its matrix, show $H\sz H=\sx$ and $H\sy H=-\sy$, and deduce $H\sx H=\sz$ again from $H^2=\Id$ without computing.

  \exercise{1} Classify each of these matrices as Hermitian, unitary, idempotent, projector, nilpotent, involutory (several or none may apply): $\begin{pmatrix}1&0\\0&0\end{pmatrix}$, $\begin{pmatrix}0&0\\1&0\end{pmatrix}$, $\tfrac12\begin{pmatrix}1&-1\\-1&1\end{pmatrix}$, $\tfrac1{\sqrt2}\begin{pmatrix}1&1\\-1&1\end{pmatrix}$, $\begin{pmatrix}1&0\\0&\ii\end{pmatrix}$, $\begin{pmatrix}2&\ii\\-\ii&3\end{pmatrix}$.

  \exercise{2} Prove: the diagonal entries of a Hermitian matrix are real; the general $2\times2$ Hermitian matrix is $\begin{pmatrix}a&c-d\ii\\c+d\ii&b\end{pmatrix}$ with $a,b,c,d\in\R$; if $A,B$ are Hermitian then $AB$ is Hermitian iff $AB=BA$, and $\ii[A,B]$ is Hermitian.

  \exercise{2} Let $P^2=P$. Show that $\Id-P$ is idempotent, that the eigenvalues of $P$ lie in $\{0,1\}$, and find a non-Hermitian idempotent $2\times2$ matrix, describing geometrically what it does.

  \exercise{3} Let $N$ be nilpotent with $N^k=0$. Show that its only eigenvalue is $0$, that $\Id-N$ is invertible with inverse $\Id+N+\dots+N^{k-1}$, and that a nilpotent Hermitian matrix is $0$.

  \exercise{2} Show that the eigenvalues of an involutory matrix are $\pm1$ and that, if $M^2=\Id$, then $P_\pm=\tfrac12(\Id\pm M)$ are idempotents with $P_++P_-=\Id$, $P_+P_-=0$ and $M=P_+-P_-$.

  \exercise{2} Show that $U$ is unitary if and only if its columns form an orthonormal basis of $\C^n$; that products and inverses of unitaries are unitary and $\abs{\det U}=1$; and that the phase gate $S=\diag(1,\ii)$ and $T=\diag(1,e^{\ii\pi/4})$ are unitary with $T^2=S$ and $S^2=\sz$.

  \exercise{1} Compute all nine products $\sigma_j\sigma_k$ and check them against Proposition~\ref{prop:pauli}. Then compute $\sz\sy\sx$ and $(\sx\sy)^2$ using only the rules of the proposition.

  \exercise{2} Define $[A,B]=AB-BA$ and $\{A,B\}=AB+BA$. Show that $[\sigma_j,\sigma_k]=2\ii\sum_l\varepsilon_{jkl}\sigma_l$ and $\{\sigma_j,\sigma_k\}=2\delta_{jk}\Id$, and prove the Jacobi identity $[A,[B,C]]+[B,[C,A]]+[C,[A,B]]=0$ for any matrices.

  \exercise{2} For $a,b\in\R^3$ write $a\cdot\sigma=a_x\sx+a_y\sy+a_z\sz$. Prove $(a\cdot\sigma)(b\cdot\sigma)=(a\cdot b)\Id+\ii(a\times b)\cdot\sigma$. Deduce that for a unit vector $n$ one has $(n\cdot\sigma)^2=\Id$, and that its eigenvalues are $\pm1$.

  \exercise{2} Decompose in the Pauli basis: $\begin{pmatrix}1&3\ii\\-\ii&1\end{pmatrix}$, $\ket0\!\bra0$, $H$ and $\begin{pmatrix}a&b\\c&d\end{pmatrix}$.
\end{exerciseset}

%======================================================================
\section{Hilbert spaces and Dirac notation}\label{sec:hilbert}
%======================================================================

\begin{objectives}
  You will compute inner products and norms, prove and use the Cauchy--Schwarz
  inequality, read and write Dirac notation, expand a vector in an orthonormal
  basis and interpret the squared coefficients, and build orthonormal bases with the
  Gram--Schmidt procedure.
\end{objectives}

A vector space has no notion of length or angle: we can add vectors, but we
cannot say whether two of them are close, or perpendicular. Quantum mechanics
needs both. The overlap between two states will be a complex number, its modulus
squared will be a probability, and ``perpendicular'' will mean ``perfectly
distinguishable''. The structure that provides them is the inner product.

\subsection{Inner products and norms}

\begin{definition}[Inner product]\label{def:inner}
  An \term{inner product} on a complex vector space $V$ is a map
  $\langle\cdot,\cdot\rangle:V\times V\to\C$ such that, for all $u,v,w\in V$ and $\lambda\in\C$:
  \begin{enumerate}
    \item[(i)] $\langle v,w\rangle=\overline{\langle w,v\rangle}$ (conjugate symmetry);
    \item[(ii)] $\langle v,v\rangle\ge0$, with $\langle v,v\rangle=0$ only for $v=0$ (positivity);
    \item[(iii)] $\langle u,v+\lambda w\rangle=\langle u,v\rangle+\lambda\langle u,w\rangle$ (linearity in the second argument).
  \end{enumerate}
  A vector space with an inner product is called an \term{inner product space}; in finite dimension, it is also a \term{Hilbert space}.
  The \term{norm} of $v$ is $\norm{v}=\sqrt{\langle v,v\rangle}$.
\end{definition}

\begin{remark}
  In infinite dimension, a Hilbert space is required to be \emph{complete}, which
  means that Cauchy sequences converge. Every finite-dimensional inner product space is
  complete automatically, and we work in finite dimension almost everywhere, so
  we will not mention completeness again.
\end{remark}

Axioms (i) and (iii) imply that the inner product is \emph{conjugate-linear} in the first argument:
$\langle\lambda v,w\rangle=\bar\lambda\langle v,w\rangle$. This is the convention of physicists; mathematicians often
put the conjugate on the second argument instead.

\begin{pitfall}
  Scalars come out of the first slot \emph{conjugated} and out of the second one unchanged.
  Forgetting this is the most common source of sign and phase errors.
\end{pitfall}

\paragraph{Examples.}
On $\C^n$ the standard inner product is $\langle v,w\rangle=\sum_i\overline{v_i}\,w_i$, so $\norm{v}^2=\sum_i\abs{v_i}^2$. On the
continuous functions $[0,1]\to\C$ we use $\langle f,g\rangle=\int_0^1\overline{f(x)}\,g(x)\dd x$. On $n\times n$ matrices,
$\langle A,B\rangle=\Tr(A^\dagger B)$, which is the standard inner product of the lists of entries.

\begin{example}\label{ex:isinner}
  Which of the following are inner products on $\C^2$? (a) $\bar v_1w_1+2\bar v_2w_2$; (b) $\bar v_1w_1-\bar v_2w_2$;
  (c) $v_1w_1+v_2w_2$; (d) $\overline{(v_1+v_2)}(w_1+w_2)$.
\end{example}
\begin{solution}
  (a) Yes. Conjugate symmetry and linearity are clear, and $f(v,v)=\abs{v_1}^2+2\abs{v_2}^2>0$ unless $v=0$.

  (b) No: positivity fails. For $v=(0,1)$ we get $f(v,v)=-1<0$.

  (c) No: there is no conjugate, so it is not conjugate symmetric. It also violates positivity:
  $f(v,v)=1+\ii^2=0$ for $v=(1,\ii)\neq0$.

  (d) No: $f(v,v)=\abs{v_1+v_2}^2$ vanishes for the non-zero vector $v=(1,-1)$.
\end{solution}

\begin{example}\label{ex:fn-inner}
  In the space of continuous functions $[0,1]\to\C$ compute $\langle 1,x\rangle$, $\norm1$ and $\norm x$.
\end{example}
\begin{solution}
  $\langle1,x\rangle=\int_0^1x\dd x=\tfrac12$, $\ \norm1^2=\int_0^11\dd x=1$, and $\norm x^2=\int_0^1x^2\dd x=\tfrac13$, so
  $\norm1=1$ and $\norm x=1/\sqrt3$. Note that $\tfrac12\le1\cdot\tfrac1{\sqrt3}\approx0.577$, which is no accident: it
  is the Cauchy--Schwarz inequality below.
\end{solution}

The next result is the most important inequality in the subject. It says that the
overlap of two vectors can never be larger than the product of their lengths.

\begin{theorem}[Cauchy--Schwarz]\label{thm:cs}
  For all $u,v$ in an inner product space, $\ \abs{\langle u,v\rangle}\le\norm u\,\norm v$, with
  equality if and only if $u$ and $v$ are linearly dependent.
\end{theorem}
\begin{proof}
  If $v=0$ both sides are $0$, so assume $v\ne0$. Split $u$ into a part along $v$ and a remainder
  (Figure~\ref{fig:cs}):
  \[
    u=\frac{\langle v,u\rangle}{\norm v^2}\,v+w,\qquad
    w:=u-\frac{\langle v,u\rangle}{\norm v^2}\,v .
  \]
  Then $w\perp v$, because $\langle v,w\rangle=\langle v,u\rangle-\frac{\langle v,u\rangle}{\norm v^2}\langle v,v\rangle=0$.
  Since the two pieces are orthogonal, the Pythagorean theorem gives
  \[
    \norm u^2=\frac{\abs{\langle v,u\rangle}^2}{\norm v^2}+\norm w^2\ \ge\ \frac{\abs{\langle v,u\rangle}^2}{\norm v^2}.
  \]
  Multiplying by $\norm v^2$ and taking square roots gives the inequality, using $\abs{\langle v,u\rangle}=\abs{\langle u,v\rangle}$.
  Equality holds exactly when $w=0$, that is, when $u$ is a multiple of $v$.
\end{proof}

\begin{figure}[htb]
  \centering
  \begin{tikzpicture}[>=stealth,scale=1.1]
    \draw[very thick,navy,->] (0,0) -- (4,0) node[right] {$v$};
    \draw[very thick,navy,->] (0,0) -- (2.6,1.8) node[above right] {$u$};
    \draw[very thick,figacc,->] (0,0) -- (2.6,0) node[midway,below] {$\dfrac{\langle v,u\rangle}{\norm v^2}\,v$};
    \draw[very thick,ink,->] (2.6,0) -- (2.6,1.8) node[midway,right] {$w$};
    \draw[ink] (2.4,0) -- (2.4,0.2) -- (2.6,0.2);
  \end{tikzpicture}
  \caption{Proof of Cauchy--Schwarz: $u$ splits into a component along $v$ and an orthogonal remainder $w$.
  The component cannot be longer than $u$ itself.}
  \label{fig:cs}
\end{figure}

\begin{corollary}[Triangle inequality]\label{cor:tri-norm}
  $\norm{u+v}\le\norm u+\norm v$. Hence $\norm{\cdot}$ is a genuine norm.
\end{corollary}
\begin{proof}
  Expanding, $\norm{u+v}^2=\norm u^2+\norm v^2+2\Re\langle u,v\rangle\le\norm u^2+\norm v^2+2\norm u\norm v=(\norm u+\norm v)^2$,
  using $\Re z\le\abs z$ and Cauchy--Schwarz. This is the same proof as in Theorem~\ref{thm:triangle}.
\end{proof}

For unit vectors, Cauchy--Schwarz says $0\le\abs{\langle u,v\rangle}\le1$: overlaps of
normalised states are always legitimate probability amplitudes. In a real space it
also lets us define the angle $\theta$ between two vectors by
$\cos\theta=\langle u,v\rangle/(\norm u\norm v)$, since the right-hand side lies in $[-1,1]$.

\subsection{Dirac notation}

Physicists write vectors in a notation due to Dirac that is compact, a little strange at first, and indispensable once you get used to it.

\begin{itemize}
  \item A vector is written as a \term{ket} $\ket\psi$. A ket is a column vector in a chosen basis.
  \item The corresponding \term{bra} $\bra\psi$ is the row vector $\ket\psi^\dagger$, the conjugate transpose.
  \item The \term{bracket} $\braket\varphi\psi$ is the inner product $\langle\varphi,\psi\rangle$: the bra $\bra\varphi$ eats the ket $\ket\psi$ and returns a number.
  \item The \term{outer product} $\ket\psi\!\bra\varphi$ is the operator $\ket\chi\mapsto\braket\varphi\chi\,\ket\psi$: a column times a row, a matrix.
\end{itemize}
In $\C^n$ this is nothing but matrix multiplication: $\braket\varphi\psi=\varphi^\dagger\psi$ is a row times a column, a $1\times1$ matrix,
while $\ket\psi\!\bra\varphi=\psi\varphi^\dagger$ is a column times a row, an $n\times n$ matrix. The rules are the ones inherited from
the inner product: $\overline{\braket\varphi\psi}=\braket\psi\varphi$, the bra of $\lambda\ket\psi$ is $\bar\lambda\bra\psi$, and $\norm{\ket\psi}^2=\braket\psi\psi$.
Throughout the book, \emph{a state is a normalised ket}, $\braket\psi\psi=1$.

A bra can be seen as a \emph{linear form}: the map $\ket w\mapsto\braket vw$ is linear, so $\bra v$ belongs to the space $V^*=L(V,\C)$ of
linear forms on $V$, the \term{dual space}. Every linear form on a finite-dimensional Hilbert space arises this way, and the correspondence
$\ket v\mapsto\bra v$ is conjugate-linear and one-to-one. That is why bras and kets are two faces of the same information.

\begin{example}\label{ex:kets}
  Let $\ket\psi=\frac1{\sqrt3}\ket0+\frac{1+\ii}{\sqrt3}\ket1$. Check that it is normalised, write $\bra\psi$, compute $\braket0\psi$,
  and find the matrix $\ket\psi\!\bra\psi$ with its trace.
\end{example}
\begin{solution}
  As a column, $\ket\psi=\frac1{\sqrt3}\begin{pmatrix}1\\1+\ii\end{pmatrix}$, and $\braket\psi\psi=\frac13(1+\abs{1+\ii}^2)=\frac13(1+2)=1$.
  The bra is the conjugate transpose, $\bra\psi=\frac1{\sqrt3}\begin{pmatrix}1&1-\ii\end{pmatrix}$, and
  $\braket0\psi=\frac1{\sqrt3}$. The outer product is the column times the row:
  \[
    \ket\psi\!\bra\psi=\frac13\begin{pmatrix}1\\1+\ii\end{pmatrix}\begin{pmatrix}1&1-\ii\end{pmatrix}
    =\frac13\begin{pmatrix}1&1-\ii\\1+\ii&2\end{pmatrix}.
  \]
  Its trace is $\tfrac13(1+2)=1$. The matrix is Hermitian, as every $\ket\psi\!\bra\psi$ is.
\end{solution}

\begin{example}\label{ex:diracsimplify}
  Simplify $\braket01$, $\braket+-$, $(\ket0\!\bra1)(\ket1\!\bra0)$, $(\ket0\!\bra1)^2$ and $(\ket+\!\bra+)\ket0$.
\end{example}
\begin{solution}
  $\braket01=0$ and $\braket+-=\tfrac12(\braket0 0-\braket01+\braket10-\braket11)=\tfrac12(1-0+0-1)=0$.
  Next, bras and kets in the middle combine: $(\ket0\!\bra1)(\ket1\!\bra0)=\ket0\braket11\bra0=\ket0\!\bra0$,
  while $(\ket0\!\bra1)^2=\ket0\braket10\bra1=0$. Finally $(\ket+\!\bra+)\ket0=\ket+\braket+0=\tfrac1{\sqrt2}\ket+$.
\end{solution}

\subsection{Orthonormal bases}

\begin{definition}[Orthogonality]\label{def:ortho}
  Two vectors $\ket u,\ket v$ are \term{orthogonal} if $\braket uv=0$. A family $\{\ket{e_i}\}$ is
  \term{orthonormal} if $\braket{e_i}{e_j}=\delta_{ij}$, which means $1$ if $i=j$ and $0$ otherwise.
  An orthonormal family that is also a basis is an \term{orthonormal basis} (ONB).
\end{definition}

Non-zero pairwise orthogonal vectors are automatically independent: if $\sum_i\lambda_i\ket{e_i}=0$, apply $\bra{e_j}$ to get
$\lambda_j\norm{e_j}^2=0$, so $\lambda_j=0$. Hence $n$ orthonormal vectors in an $n$-dimensional space are an ONB. The point of an ONB
is that coordinates become trivial to compute.

\begin{theorem}[Expansion and Parseval]\label{thm:parseval}
  Let $\{\ket{e_i}\}$ be an ONB of $\Hil$. For every $\ket\psi\in\Hil$,
  \[
    \ket\psi=\sum_i\braket{e_i}\psi\,\ket{e_i},\qquad
    \norm{\psi}^2=\sum_i\abs{\braket{e_i}\psi}^2 .
  \]
\end{theorem}
\begin{proof}
  Write $\ket\psi=\sum_j\lambda_j\ket{e_j}$. Applying $\bra{e_i}$ gives $\braket{e_i}\psi=\sum_j\lambda_j\delta_{ij}=\lambda_i$, which is the first claim.
  For the second, $\norm\psi^2=\braket\psi\psi=\sum_{i,j}\overline{\lambda_i}\lambda_j\braket{e_i}{e_j}=\sum_i\abs{\lambda_i}^2$.
\end{proof}

No linear system has to be solved: the coordinate along $\ket{e_i}$ is the overlap $\braket{e_i}\psi$. For a normalised state, Parseval
says that the numbers $\abs{\braket{e_i}\psi}^2$ add up to $1$; when the postulates of Chapter~\ref{ch:postulates} come into play, they will be
the probabilities of the outcomes of a measurement in that basis. Reading the first formula as $\ket\psi=\bigl(\sum_i\ket{e_i}\!\bra{e_i}\bigr)\ket\psi$ gives the
\term{completeness relation}
\[
  \sum_i\ket{e_i}\!\bra{e_i}=\Id,
\]
valid for every ONB, which lets us ``insert an identity'' anywhere in a computation. For instance, inserting it on both sides of an operator $A$ gives
\[
  A=\sum_{i,j}\ket{e_i}\bra{e_i}A\ket{e_j}\bra{e_j}=\sum_{i,j}A_{ij}\ket{e_i}\!\bra{e_j},\qquad A_{ij}=\bra{e_i}A\ket{e_j},
\]
so the matrix element $A_{ij}$ is the $i$th coordinate of $A\ket{e_j}$, in agreement with the rule ``columns are images'' of Section~\ref{sec:vector-spaces}.

\begin{example}\label{ex:parseval}
  Expand $\ket\psi=\frac1{\sqrt3}\bigl(\ket0+(1+\ii)\ket1\bigr)$ in the ONBs $\{\ket0,\ket1\}$, $\{\ket+,\ket-\}$ and $\{\ket{+\ii},\ket{-\ii}\}$, and check Parseval.
\end{example}
\begin{solution}
  \emph{Computational basis.} $\braket0\psi=\frac1{\sqrt3}$ and $\braket1\psi=\frac{1+\ii}{\sqrt3}$. Squared moduli: $\tfrac13$ and $\tfrac23$, summing to $1$.

  \emph{Basis $\ket\pm$.} The bras are $\bra\pm=\frac1{\sqrt2}\begin{pmatrix}1&\pm1\end{pmatrix}$, so
  \[
    \braket+\psi=\frac{1}{\sqrt6}\bigl(1+(1+\ii)\bigr)=\frac{2+\ii}{\sqrt6},\qquad
    \braket-\psi=\frac{1}{\sqrt6}\bigl(1-(1+\ii)\bigr)=\frac{-\ii}{\sqrt6}.
  \]
  Squared moduli: $\tfrac56$ and $\tfrac16$, summing to $1$.

  \emph{Basis $\ket{\pm\ii}$.} Here $\bra{\pm\ii}=\frac1{\sqrt2}\begin{pmatrix}1&\mp\ii\end{pmatrix}$ (the second entry is conjugated), so
  \[
    \braket{+\ii}\psi=\frac1{\sqrt6}\bigl(1-\ii(1+\ii)\bigr)=\frac{2-\ii}{\sqrt6},\qquad
    \braket{-\ii}\psi=\frac1{\sqrt6}\bigl(1+\ii(1+\ii)\bigr)=\frac{\ii}{\sqrt6},
  \]
  with squared moduli $\tfrac56$ and $\tfrac16$, again summing to $1$. The same state has three different sets of coordinates, one for each point of view.
\end{solution}

\begin{pitfall}
  The formulas of Theorem~\ref{thm:parseval} need an \emph{orthonormal} basis. In $\R^2$, for the independent but
  non-orthogonal basis $(1,0),(1,1)$, the vector $(3,5)$ is $-2(1,0)+5(1,1)$, whereas the overlaps with the basis vectors are $3$ and $8$.
\end{pitfall}

\subsection{Gram--Schmidt}

Do orthonormal bases always exist? Yes, and the proof is a recipe. Given any basis, we make it orthonormal one vector at a time,
subtracting at each step the components along the vectors already produced (Figure~\ref{fig:gs}).

\begin{theorem}[Gram--Schmidt]\label{thm:gs}
  Let $\ket{v_1},\dots,\ket{v_n}$ be a basis of $\Hil$. Define recursively
  \[
    \ket{u_k}=\ket{v_k}-\sum_{j<k}\braket{e_j}{v_k}\ket{e_j},\qquad \ket{e_k}=\frac{\ket{u_k}}{\norm{u_k}},\qquad k=1,\dots,n .
  \]
  Then $\{\ket{e_k}\}$ is an orthonormal basis, and $\Span\{e_1,\dots,e_m\}=\Span\{v_1,\dots,v_m\}$ for every $m$.
\end{theorem}
\begin{proof}
  We argue by induction on $k$. Suppose $e_1,\dots,e_{k-1}$ are orthonormal and span the same space as $v_1,\dots,v_{k-1}$.
  For $i<k$, $\braket{e_i}{u_k}=\braket{e_i}{v_k}-\sum_{j<k}\braket{e_j}{v_k}\delta_{ij}=0$, so $u_k$ is orthogonal to all previous $e_i$.
  Moreover $u_k\neq0$: otherwise $v_k$ would be a combination of $e_1,\dots,e_{k-1}$, hence of $v_1,\dots,v_{k-1}$, contradicting independence.
  So $e_k$ is well defined, of norm $1$, and $\Span\{e_1,\dots,e_k\}=\Span\{v_1,\dots,v_k\}$.
\end{proof}

In particular, every finite-dimensional Hilbert space has an orthonormal basis.

\begin{figure}[htb]
  \centering
  \begin{tikzpicture}[>=stealth,scale=1.1]
    \draw[very thick,navy,->] (0,0) -- (3,0) node[right] {$e_1=v_1/\norm{v_1}$};
    \draw[very thick,navy!60,->] (0,0) -- (1.6,2.0) node[above] {$v_2$};
    \draw[dashed,ink] (1.6,2.0) -- (1.6,0);
    \draw[very thick,figacc,->] (0,0) -- (1.6,0) node[midway,below] {$\langle e_1,v_2\rangle e_1$};
    \draw[very thick,figacc,->] (1.6,0) -- (1.6,2.0) node[midway,right] {$u_2$};
    \draw[ink] (1.4,0) -- (1.4,0.2) -- (1.6,0.2);
  \end{tikzpicture}
  \caption{Gram--Schmidt, second step: removing from $v_2$ its component along $e_1$ leaves $u_2\perp e_1$,
  which after normalisation is $e_2$.}
  \label{fig:gs}
\end{figure}

\begin{example}\label{ex:gs-c2}
  Apply Gram--Schmidt to $\ket{v_1}=\ket0+\ket1$, $\ket{v_2}=\ket0$ in $\C^2$, and then in the reverse order.
\end{example}
\begin{solution}
  Since $\norm{v_1}=\sqrt2$, $\ket{e_1}=\frac1{\sqrt2}(\ket0+\ket1)=\ket+$. Then $\braket{e_1}{v_2}=\braket+0=\tfrac1{\sqrt2}$, and
  \[
    \ket{u_2}=\ket0-\tfrac1{\sqrt2}\ket+=\ket0-\tfrac12(\ket0+\ket1)=\tfrac12(\ket0-\ket1),\qquad
    \ket{e_2}=\frac{\ket{u_2}}{\norm{u_2}}=\ket-,
  \]
  because $\norm{u_2}=\tfrac1{\sqrt2}$. The result is the basis $\{\ket+,\ket-\}$.

  In the reverse order, $\ket{e_1}=\ket0$ and $\ket{u_2}=(\ket0+\ket1)-\braket0{v_1}\ket0=\ket1$, so $\ket{e_2}=\ket1$:
  now we obtain $\{\ket0,\ket1\}$. \emph{The outcome depends on the order of the input.}
\end{solution}

\begin{example}\label{ex:gs-r3}
  Apply Gram--Schmidt to $v_1=(1,1,0)$, $v_2=(1,0,1)$, $v_3=(0,1,1)$ in $\R^3$.
\end{example}
\begin{solution}
  $e_1=\frac1{\sqrt2}(1,1,0)$. Since $\langle e_1,v_2\rangle=\tfrac1{\sqrt2}$, we get
  $u_2=v_2-\tfrac1{\sqrt2}e_1=(1,0,1)-(\tfrac12,\tfrac12,0)=(\tfrac12,-\tfrac12,1)$, with $\norm{u_2}^2=\tfrac14+\tfrac14+1=\tfrac32$, so
  $e_2=\frac1{\sqrt6}(1,-1,2)$.
  For the third vector, $\langle e_1,v_3\rangle=\tfrac1{\sqrt2}$ and $\langle e_2,v_3\rangle=\tfrac1{\sqrt6}$, hence
  \[
    u_3=(0,1,1)-(\tfrac12,\tfrac12,0)-\tfrac16(1,-1,2)=(-\tfrac23,\tfrac23,\tfrac23),\qquad
    e_3=\tfrac1{\sqrt3}(-1,1,1).
  \]
  \emph{Check:} $e_1\cdot e_3=0$, $e_2\cdot e_3=\tfrac1{\sqrt{18}}(-1-1+2)=0$.
\end{solution}

\begin{exerciseset}
  \exercise{1} From axioms (i)--(iii) prove that $\langle\lambda v,w\rangle=\bar\lambda\langle v,w\rangle$, that $\langle v+v',w\rangle=\langle v,w\rangle+\langle v',w\rangle$, that $\langle0,w\rangle=0$,
  and that if $\langle v,w\rangle=0$ for all $w$ then $v=0$.

  \exercise{1} Decide whether each is an inner product on $\C^2$: $\bar v_1w_1+\bar v_2w_2$; $\bar v_1w_1-\bar v_2w_2$; $\bar v_1w_1$; $\bar v_1w_1+\bar v_2w_2+\bar v_1w_2+\bar v_2w_1$.

  \exercise{2} Let $A\in M_n(\C)$ and $f(v,w)=v^\dagger Aw$. Show that axiom (i) holds iff $A$ is Hermitian, and that $f$ is an inner product iff
  $A$ is Hermitian and $v^\dagger Av>0$ for all $v\neq0$. For which $c\in\C$ is $A=\begin{pmatrix}1&c\\\bar c&1\end{pmatrix}$ an inner product matrix?

  \exercise{2} On $C([0,1],\C)$ let $\langle f,g\rangle=\int_0^1\overline f\,g\dd x$. Check that this is an inner product (where do you use continuity?),
  compute $\langle1,x\rangle$ and $\norm x$, and show that $\{e^{2\pi\ii nx}\}_{n\in\Z}$ is an orthonormal system.

  \exercise{2} On $M_n(\C)$ define $\langle A,B\rangle=\Tr(A^\dagger B)$. Show that it is an inner product with $\norm A^2=\sum_{ij}\abs{A_{ij}}^2$,
  and that $\{\tfrac1{\sqrt2}\sigma_\mu\}_{\mu=0}^3$ is an orthonormal basis of $M_2(\C)$.

  \exercise{2} Show that $\norm{x}_1=\sum_i\abs{x_i}$ and $\norm x_\infty=\max_i\abs{x_i}$ are norms on $\C^n$, and sketch the unit balls of $\norm\cdot_1$, $\norm\cdot_2$,
  $\norm\cdot_\infty$ in $\R^2$.

  \exercise{3} Prove the polarisation identity, $\langle u,v\rangle=\tfrac14\bigl(\norm{u+v}^2-\norm{u-v}^2-\ii\norm{u+\ii v}^2+\ii\norm{u-\ii v}^2\bigr)$, with our convention. Conclude that the norm determines the inner
  product. Prove also the parallelogram law $\norm{u+v}^2+\norm{u-v}^2=2\norm u^2+2\norm v^2$ and show that $\norm\cdot_1$ and $\norm\cdot_\infty$ on $\R^2$ violate it.

  \exercise{1} Let $\ket\psi=\frac1{\sqrt3}\ket0+\frac{1+\ii}{\sqrt3}\ket1$ and $\ket\varphi=\frac1{\sqrt2}(\ket0-\ii\ket1)$. Write them as columns and the bras as rows, compute $\braket\psi\psi$, $\braket\varphi\varphi$, $\braket\varphi\psi$,
  and check $\braket\varphi\psi=\overline{\braket\psi\varphi}$. Compute $\ket\psi\!\bra\varphi$, $\ket\psi\!\bra\psi$ and $\ket\varphi\!\bra\varphi$: which are Hermitian, and what are their ranks and traces?

  \exercise{1} Simplify: $\bra0\sx\ket1$, $\bra1\sy\ket0$, $\bra+\sz\ket-$, $(\ket0\!\bra1)(\ket1\!\bra0)$, $(\ket+\!\bra+)\ket0$, $\bra{+\ii}\sy\ket{+\ii}$.

  \exercise{2} Prove: $(\ket\psi\!\bra\varphi)^\dagger=\ket\varphi\!\bra\psi$; $(\ket\psi\!\bra\varphi)(\ket\chi\!\bra\xi)=\braket\varphi\chi\,\ket\psi\!\bra\xi$; and $(\lambda\ket\psi)^\dagger=\bar\lambda\bra\psi$.

  \exercise{1} Write $\sx,\sy,\sz,H$ as $\sum_{ij}A_{ij}\ket i\!\bra j$, and compute the entries $A_{ij}=\bra{i}A\ket{j}$ of $\sy$ in the basis $\{\ket+,\ket-\}$.

  \exercise{1} Check that $\sum_i\ket{e_i}\!\bra{e_i}=\Id$ for the three bases $\{\ket0,\ket1\}$, $\{\ket\pm\}$, $\{\ket{\pm\ii}\}$.

  \exercise{2} For $\ket\psi=\cos\frac\theta2\ket0+e^{\ii\phi}\sin\frac\theta2\ket1$, compute $\bra\psi\sx\ket\psi$, $\bra\psi\sy\ket\psi$ and $\bra\psi\sz\ket\psi$, and show that the squares add up to $1$.

  \exercise{2} Show that $\ket\psi$ and $e^{\ii\gamma}\ket\psi$ give the same value of $\bra\psi A\ket\psi$ for every operator $A$ and of $\abs{\braket\varphi\psi}^2$ for every $\varphi$. Then show that
  $\frac1{\sqrt2}(\ket0+\ket1)$ and $\frac1{\sqrt2}(\ket0-\ket1)$ are \emph{not} equivalent in this sense, by finding an operator that tells them apart.

  \exercise{1} Which pairs are orthogonal? $\ket0,\ket+$; $\ket+,\ket-$; $\ket{+\ii},\ket{-\ii}$; $(1,\ii)$ and $(\ii,1)$; $(1,\ii)$ and $(1,-\ii)$; $(1,1,\ii)$ and $(1,-1,0)$.

  \exercise{2} Find all unit vectors of $\C^2$ orthogonal to $\frac1{\sqrt3}\bigl(\ket0+\sqrt2\,e^{\ii\pi/3}\ket1\bigr)$.

  \exercise{2} Let $\{\ket{e_i}\}$ be an ONB. Prove $\ket\psi=\sum_i\braket{e_i}\psi\ket{e_i}$ and $\norm\psi^2=\sum_i\abs{\braket{e_i}\psi}^2$ (Parseval), and show with an example in $\R^2$ that both fail for a non-orthogonal basis.

  \exercise{2} Expand $\ket\psi=\frac1{\sqrt3}\ket0+\frac{1+\ii}{\sqrt3}\ket1$ in the three bases of Example~\ref{ex:parseval} (checking the computations there), and interpret the squared coefficients.

  \exercise{2} Prove Pythagoras: $\braket uv=0\Rightarrow\norm{u+v}^2=\norm u^2+\norm v^2$. Show that the converse holds over $\R$ but fails over $\C$ (try $v=\ii u$).

  \exercise{3} Prove Bessel's inequality: if $\{\ket{e_1},\dots,\ket{e_k}\}$ is orthonormal but not necessarily a basis, then $\sum_{i\le k}\abs{\braket{e_i}\psi}^2\le\norm\psi^2$, with equality iff $\ket\psi\in\Span\{e_i\}$.

  \exercise{3} Two ONBs $\{\ket{e_i}\},\{\ket{f_j}\}$ of $\C^n$ are \emph{mutually unbiased} if $\abs{\braket{e_i}{f_j}}^2=1/n$ for all $i,j$. Show that the eigenbases of $\sx,\sy,\sz$ are pairwise
  mutually unbiased. Interpret: a state from one basis gives completely random outcomes when measured in another.

  \exercise{2} Let $B=\{(1,0),(1,1)\}\subset\R^2$. Find the dual basis $\{v_1^*,v_2^*\}$ as row vectors (so that $v_i^*(v_j)=\delta_{ij}$) and use it to compute the coordinates of $(3,5)$ in $B$.

  \exercise{3} Show that the dual basis of $\{\ket0,\ket+\}$ in $\C^2$ does not consist of the bras $\bra0,\bra+$, and explain why for an orthonormal basis it does.

  \exercise{1} Redo Example~\ref{ex:gs-c2} in both orders. Then apply Gram--Schmidt to $(1,1,0),(1,0,1),(0,1,1)$ in $\R^3$, and to $(1,\ii,0),(0,1,1),(1,0,1)$ in $\C^3$ in both orders.

  \exercise{2} Apply Gram--Schmidt to $(1,\ii),(\ii,1)$ in $\C^2$. What goes wrong with $(1,\ii),(\ii,-1)$, and why?

  \exercise{2} Prove that $\ket{u_k}\perp\ket{e_j}$ for $j<k$, that $u_k\neq0$, and that $\Span\{e_1,\dots,e_m\}=\Span\{v_1,\dots,v_m\}$. Show that if the $v_i$ are already orthogonal, Gram--Schmidt only normalises them.

  \exercise{3} Apply Gram--Schmidt to $\{1,x,x^2,x^3\}$ with $\langle f,g\rangle=\int_{-1}^1fg\dd x$. Compare with the Legendre polynomials $1,\ x,\ \tfrac12(3x^2-1),\ \tfrac12(5x^3-3x)$.

  \exercise{2} Extend $\ket{e_1}=\tfrac1{\sqrt3}(1,1,1)$ to an orthonormal basis of $\C^3$ using Gram--Schmidt. Find a different completion using $\tfrac1{\sqrt3}(1,\omega,\omega^2)$ and $\tfrac1{\sqrt3}(1,\omega^2,\omega^4)$ with $\omega=e^{2\pi\ii/3}$.
\end{exerciseset}

%======================================================================
\section{Operators and their spectra}\label{sec:operators}
%======================================================================

\begin{objectives}
  You will compute adjoints and recognise Hermitian and unitary operators; find eigenvalues and eigenvectors; diagonalise Hermitian and unitary operators with the spectral theorem; use projectors; and apply functions, such as the exponential, to operators.
\end{objectives}

\begin{recall}
  From Section~\ref{sec:hilbert}: inner products, norms, orthonormal bases, Dirac notation and the completeness relation $\sum_i\ket{e_i}\!\bra{e_i}=\Id$. From Section~\ref{sec:vector-spaces}: matrices of linear maps, the conjugate transpose $M^\dagger$, the trace and the determinant, and the Pauli matrices.
\end{recall}

Operators are the active ingredients of quantum mechanics. Each quantity that we can measure will be a Hermitian operator, each reversible change of a quantum system will be a unitary one, and in both cases the interesting information is carried by eigenvalues and eigenvectors. An \term{operator} on $\Hil$ is just a linear map $A\in L(\Hil)$; in an ONB it is a square matrix. From now on $\Hil$ is a finite-dimensional Hilbert space.

\subsection{The adjoint}

In Section~\ref{sec:vector-spaces} we defined the adjoint of a matrix by a recipe, ``transpose and conjugate''. A recipe tells us how to compute but not what the object is. The inner product gives a recipe-free definition.

\begin{definition}[Adjoint]\label{def:adjoint}
  The \term{adjoint} of $A\in L(\Hil)$ is the unique operator $A^\dagger\in L(\Hil)$ such that
  \[
    \braket{\varphi}{A\psi}=\braket{A^\dagger\varphi}{\psi}\qquad\text{for all }\ket\varphi,\ket\psi\in\Hil .
  \]
\end{definition}

In words: an operator may jump to the other side of the bracket, but it must put on its dagger. Uniqueness is easy (if $B$ and $B'$ both work then $\langle(B-B')\varphi,\psi\rangle=0$ for all $\psi$, so $(B-B')\varphi=0$), and existence follows from the next theorem, which also recovers the recipe.

\begin{theorem}[Recovering the recipe]\label{thm:adjmatrix}
  Let $\{\ket{e_i}\}$ be an \emph{orthonormal} basis and $A_{ij}=\bra{e_i}A\ket{e_j}$. Then $(A^\dagger)_{ij}=\overline{A_{ji}}$.
\end{theorem}
\begin{proof}
  \begin{steps}
    \item \emph{Definition of the matrix element.} $(A^\dagger)_{ij}=\bra{e_i}A^\dagger\ket{e_j}$.
    \item \emph{Conjugate symmetry.} $=\overline{\braket{A^\dagger e_j}{e_i}}$.
    \item \emph{Definition of the adjoint.} $=\overline{\braket{e_j}{Ae_i}}$.
    \item \emph{Matrix element again.} $=\overline{A_{ji}}$. \qedhere
  \end{steps}
\end{proof}

\begin{pitfall}
  ``Transpose and conjugate'' is only valid in an \emph{orthonormal} basis. In a non-orthonormal basis the matrix of $A^\dagger$ is a different matrix.
\end{pitfall}

\begin{example}\label{ex:adjoint}
  Find the adjoint of $A=\begin{pmatrix}1&\ii\\2&3-\ii\end{pmatrix}$, and of $\ket0\!\bra1+\ii\ket1\!\bra1$.
\end{example}
\begin{solution}
  \begin{steps}
    \item \emph{Off-diagonal entries.} The $(1,2)$ entry of $A^\dagger$ is $\overline{A_{21}}=2$ and the $(2,1)$ entry is $\overline{A_{12}}=-\ii$.
    \item \emph{Diagonal entries} are conjugated: $1$ and $\overline{3-\ii}=3+\ii$.
    \item \emph{Result.} $A^\dagger=\begin{pmatrix}1&2\\-\ii&3+\ii\end{pmatrix}$.
    \item \emph{Dirac form.} Use $(\ket\psi\!\bra\varphi)^\dagger=\ket\varphi\!\bra\psi$ and conjugate scalars: $(\ket0\!\bra1+\ii\ket1\!\bra1)^\dagger=\ket1\!\bra0-\ii\ket1\!\bra1$.
  \end{steps}
\end{solution}

The adjoint obeys a handful of rules, all easy consequences of the definition:
\[
  (A^\dagger)^\dagger=A,\quad (A+B)^\dagger=A^\dagger+B^\dagger,\quad (\lambda A)^\dagger=\bar\lambda A^\dagger,\quad (AB)^\dagger=B^\dagger A^\dagger,\quad (A^{-1})^\dagger=(A^\dagger)^{-1}.
\]
For instance, $\braket\varphi{AB\psi}=\braket{A^\dagger\varphi}{B\psi}=\braket{B^\dagger A^\dagger\varphi}{\psi}$ proves the fourth: \emph{the order flips}. In Dirac notation, $(A\ket\psi)^\dagger=\bra\psi A^\dagger$ and $\bra\varphi A\ket\psi=\overline{\bra\psi A^\dagger\ket\varphi}$.

\begin{remark}
  The definition makes sense far beyond matrices. On functions that vanish at infinity, with $\langle f,g\rangle=\int\overline f\,g\dd x$, integration by parts gives $\langle f,g'\rangle=-\langle f',g\rangle$, so $(\frac{d}{dx})^\dagger=-\frac{d}{dx}$ and the momentum operator $\hat p=-\ii\hbar\frac{d}{dx}$ equals its own adjoint, with no matrix in sight.
\end{remark}

\subsection{Hermitian and unitary operators}

The adjoint singles out the two families of operators that quantum mechanics is built on.

\begin{definition}[Hermitian and unitary operators]\label{def:herm-unit}
  An operator $A$ is \term{Hermitian} (or self-adjoint) if $A^\dagger=A$, and \term{unitary} if $U^\dagger U=UU^\dagger=\Id$.
\end{definition}

A Hermitian operator is one that moves freely across the bracket, $\braket{\varphi}{A\psi}=\braket{A\varphi}{\psi}$. It follows that its \term{expectation value} $\bra\psi A\ket\psi$ is real, since it equals its own conjugate: $\overline{\bra\psi A\ket\psi}=\bra\psi A^\dagger\ket\psi=\bra\psi A\ket\psi$. Real numbers are what an instrument displays, so Hermitian operators will be the observables.
A unitary operator preserves every inner product, $\braket{U\varphi}{U\psi}=\braket\varphi{U^\dagger U\psi}=\braket\varphi\psi$, and in particular every norm: it maps states to states, reversibly, as $U^{-1}=U^\dagger$. Unitary operators will be the quantum gates, and the next chapter makes this precise.

\begin{example}\label{ex:classify}
  Classify as Hermitian and/or unitary: $\sy$, $\ \tfrac1{\sqrt2}\begin{pmatrix}1&1\\-1&1\end{pmatrix}$, $\ \diag(1,\ii)$, $\ \begin{pmatrix}2&\ii\\-\ii&3\end{pmatrix}$.
\end{example}
\begin{solution}
  \begin{steps}
    \item \emph{$\sy$} is both: it equals its adjoint and squares to $\Id$ (Proposition~\ref{prop:pauli}).
    \item \emph{$U=\tfrac1{\sqrt2}\begin{pmatrix}1&1\\-1&1\end{pmatrix}$:} $U^\dagger=\tfrac1{\sqrt2}\begin{pmatrix}1&-1\\1&1\end{pmatrix}$ and $U^\dagger U=\tfrac12\begin{pmatrix}2&0\\0&2\end{pmatrix}=\Id$, so it is unitary (a rotation by $-45^\circ$). It is not Hermitian, as $U^\dagger\neq U$.
    \item \emph{$\diag(1,\ii)$} is unitary, since its adjoint is $\diag(1,-\ii)$ and the product with it is $\Id$, but not Hermitian.
    \item \emph{$M=\begin{pmatrix}2&\ii\\-\ii&3\end{pmatrix}$} is Hermitian, as $M^\dagger=M$. It is not unitary: $M^\dagger M=M^2=\begin{pmatrix}5&5\ii\\-5\ii&10\end{pmatrix}\neq\Id$.
  \end{steps}
\end{solution}

\subsection{Eigenvalues and eigenvectors}

\begin{definition}[Eigenvalue and eigenvector]\label{def:eigen}
  Let $A\in L(\Hil)$. A non-zero vector $\ket v$ is an \term{eigenvector} of $A$ with \term{eigenvalue} $\lambda\in\C$ if $A\ket v=\lambda\ket v$, equivalently $(A-\lambda\Id)\ket v=0$. The set of eigenvalues is the \term{spectrum} $\spec(A)$. The eigenvalues are the roots of the \term{characteristic polynomial} $p(\lambda)=\det(A-\lambda\Id)$.
\end{definition}

An eigenvector is a direction that $A$ does not turn: it only rescales it by $\lambda$. The reason for working over $\C$ is that every non-constant polynomial has a root there, so every operator has at least one eigenvalue; over $\R$ this fails. The recipe to find eigenvalues and eigenvectors never changes:
\begin{steps}
  \item compute $p(\lambda)=\det(A-\lambda\Id)$ and solve $p(\lambda)=0$;
  \item for each root $\lambda$, solve the linear system $(A-\lambda\Id)v=0$;
  \item normalise the solutions if an orthonormal basis is wanted.
\end{steps}

\begin{example}\label{ex:eigen}
  Find the eigenvalues and eigenvectors of $\sx$, of $A=\begin{pmatrix}2&1\\1&2\end{pmatrix}$, and of the rotation $R_\theta=\begin{pmatrix}\cos\theta&-\sin\theta\\\sin\theta&\cos\theta\end{pmatrix}$ for $\theta\notin\pi\Z$.
\end{example}
\begin{solution}
  \begin{steps}
    \item \emph{$\sx$: eigenvalues.} $p(\lambda)=\det\begin{pmatrix}-\lambda&1\\1&-\lambda\end{pmatrix}=\lambda^2-1$, so $\lambda=\pm1$.
    \item \emph{$\sx$: eigenvectors.} Solving $\sx\ket v=\pm\ket v$ gives $\ket v=(1,\pm1)$, that is $\ket\pm=\frac1{\sqrt2}(\ket0\pm\ket1)$. Thus $\sx\ket\pm=\pm\ket\pm$.
    \item \emph{$A$: eigenvalues.} $p(\lambda)=(2-\lambda)^2-1=(\lambda-1)(\lambda-3)$.
    \item \emph{$A$: eigenvectors.} For $\lambda=3$, $(A-3\Id)v=0$ gives $v_1=v_2$, so $v=(1,1)$, that is $\ket+$. For $\lambda=1$ we get $v=(1,-1)$, that is $\ket-$. (Indeed $A=2\Id+\sx$.)
    \item \emph{$R_\theta$: eigenvalues.} $p(\lambda)=\lambda^2-2\cos\theta\,\lambda+1$ with roots $\lambda=\cos\theta\pm\ii\sin\theta=e^{\pm\ii\theta}$, real only if $\theta\in\pi\Z$.
    \item \emph{$R_\theta$: an eigenvector.} For $\lambda=e^{\ii\theta}$ take $(1,-\ii)$: $R_\theta(1,-\ii)=(\cos\theta+\ii\sin\theta,\ \sin\theta-\ii\cos\theta)=e^{\ii\theta}(1,-\ii)$. A rotation of the plane has no real eigenvector, which is natural, and has two complex ones.
  \end{steps}
\end{solution}

\begin{example}[Eigen-decomposition of $\sy$]\label{ex:eigen-y}
  Find the eigenvalues and an orthonormal eigenbasis of $\sy$.
\end{example}
\begin{solution}
  \begin{steps}
    \item \emph{Characteristic polynomial.} $p(\lambda)=\det\begin{pmatrix}-\lambda&-\ii\\\ii&-\lambda\end{pmatrix}=\lambda^2-(-\ii)(\ii)=\lambda^2-1$, so $\lambda=\pm1$.
    \item \emph{$\lambda=+1$.} $(\sy-\Id)v=0$: $-v_1-\ii v_2=0$, so $v_1=-\ii v_2$; take $v=(1,\ii)$ (indeed $-\ii\cdot\ii=1$). Normalised: $\ket{+\ii}=\frac1{\sqrt2}(\ket0+\ii\ket1)$.
    \item \emph{$\lambda=-1$.} $(\sy+\Id)v=0$: $v_1=\ii v_2$; take $v=(1,-\ii)$, normalised $\ket{-\ii}=\frac1{\sqrt2}(\ket0-\ii\ket1)$.
    \item \emph{Orthogonality.} $\braket{+\ii}{-\ii}=\tfrac12(1+(-\ii)(-\ii))=\tfrac12(1-1)=0$, as Theorem~\ref{thm:spectrum} will promise for a Hermitian operator.
  \end{steps}
\end{solution}

The two families of the previous subsection have a very constrained spectrum, and this is what makes them physically meaningful.

\begin{theorem}\label{thm:spectrum}
  \begin{enumerate}
    \item[(a)] If $H$ is Hermitian, its eigenvalues are real, and eigenvectors for different eigenvalues are orthogonal.
    \item[(b)] If $U$ is unitary, its eigenvalues have modulus $1$: $\lambda=e^{\ii\theta}$ with $\theta\in\R$.
  \end{enumerate}
\end{theorem}
\begin{proof}
  \begin{steps}
    \item (a) \emph{Real eigenvalues.} Let $H\ket v=\lambda\ket v$ with $\norm v=1$. Then $\lambda=\bra vH\ket v$, which is real because $H$ is Hermitian.
    \item (a) \emph{Orthogonality.} If also $H\ket w=\mu\ket w$ with $\mu\neq\lambda$, then $\lambda\braket vw=\braket{Hv}{w}=\braket v{Hw}=\mu\braket vw$ (using $\lambda$ real). So $(\lambda-\mu)\braket vw=0$ and $\braket vw=0$.
    \item (b) Let $U\ket v=\lambda\ket v$ with $\norm v=1$. Since $U$ preserves norms, $1=\braket{Uv}{Uv}=\abs\lambda^2\braket vv=\abs\lambda^2$. \qedhere
  \end{steps}
\end{proof}

In words: measurement outcomes, which are eigenvalues of Hermitian operators, are real, and distinct outcomes correspond to orthogonal states, which are perfectly distinguishable. Unitary evolution cannot make eigenvalues grow or decay: it can only rotate their phase.

\begin{example}[A unitary with phase eigenvalues]\label{ex:unitary-eig}
  Find the eigenvalues and eigenvectors of $S=\diag(1,\ii)$ and of $R_{\pi/2}=\begin{pmatrix}0&-1\\1&0\end{pmatrix}$, and check Theorem~\ref{thm:spectrum}(b).
\end{example}
\begin{solution}
  \begin{steps}
    \item \emph{$S$.} It is diagonal, so the eigenvalues are $1$ and $\ii=e^{\ii\pi/2}$, with eigenvectors $\ket0,\ket1$.
    \item \emph{$R_{\pi/2}$.} By Example~\ref{ex:eigen} with $\theta=\pi/2$, the eigenvalues are $e^{\pm\ii\pi/2}=\pm\ii$, with eigenvectors $(1,\mp\ii)$.
    \item \emph{Check.} All four eigenvalues have modulus $1$. Both operators are unitary: they rotate, one the phase of $\ket1$, the other the plane by a quarter turn.
  \end{steps}
\end{solution}

\subsection{The spectral theorem}

We saw in Example~\ref{ex:changebasis} that a good choice of basis can make a map very simple. The spectral theorem says that for Hermitian and unitary operators the best basis is an orthonormal one made of eigenvectors.

\begin{theorem}[Spectral theorem]\label{thm:spectral}
  Let $H\in L(\Hil)$ be Hermitian. Then $\Hil$ has an orthonormal basis $\{\ket{v_1},\dots,\ket{v_n}\}$ of eigenvectors of $H$, with real eigenvalues $\lambda_1,\dots,\lambda_n$. In matrix form, there is a unitary $U$ with
  \[
    U^\dagger HU=D=\diag(\lambda_1,\dots,\lambda_n),\qquad H=\sum_{i=1}^n\lambda_i\ket{v_i}\!\bra{v_i}.
  \]
  The same holds for every unitary operator, with eigenvalues $e^{\ii\theta_i}$. More generally it holds for every \emph{normal} operator ($AA^\dagger=A^\dagger A$).
\end{theorem}
\begin{proof}
  We prove it for a Hermitian $H$ by induction on $n=\dim\Hil$. For $n=1$ there is nothing to prove.
  \begin{steps}
    \item \emph{An eigenvector exists.} For $n>1$ the characteristic polynomial has a root $\lambda_1$ over $\C$, so $H$ has a unit eigenvector $\ket{v_1}$.
    \item \emph{Its orthogonal complement is invariant.} Let $W=\{\ket w:\braket{v_1}w=0\}$. If $\ket w\in W$ then $\braket{v_1}{Hw}=\braket{Hv_1}{w}=\lambda_1\braket{v_1}w=0$, so $H\ket w\in W$.
    \item \emph{Induct.} The restriction of $H$ to $W$ is Hermitian and $\dim W=n-1$, so by induction $W$ has an ONB of eigenvectors of $H$.
    \item \emph{Assemble.} Adding $\ket{v_1}$ gives an ONB of $\Hil$. The matrix $U$ whose columns are the $\ket{v_i}$ is unitary, and $U^\dagger HU=D$ because $H\ket{v_j}=\lambda_j\ket{v_j}$.
    \item \emph{The decomposition.} $H=\sum_i\lambda_i\ket{v_i}\!\bra{v_i}$ is the matrix form $D$ read back through the completeness relation. \qedhere
  \end{steps}
\end{proof}

The formula $H=\sum_i\lambda_i\ket{v_i}\!\bra{v_i}$ is called the \term{spectral decomposition} of $H$. It writes $H$ as a combination of rank-one pieces, one for each direction in which $H$ acts by a number. In the right orthonormal basis, a Hermitian operator just rescales each axis independently (Figure~\ref{fig:ellipse}), a unitary one just rotates each by a phase, and a coupled problem of dimension $n$ falls apart into $n$ independent scalar problems.

\begin{figure}[htb]
  \centering
  \begin{tikzpicture}[>=stealth,scale=1.15]
    \draw[->,ink!70] (-3.4,0)--(3.4,0); \draw[->,ink!70] (0,-3.4)--(0,3.4);
    \draw[ink,thick] (0,0) circle (1);
    \begin{scope}[rotate=45]
      \draw[figacc,very thick] (0,0) ellipse (3 and 1);
      \draw[dashed,ink!60] (-3.3,0)--(3.3,0);
      \draw[dashed,ink!60] (0,-1.5)--(0,1.5);
      \draw[very thick,navy,->] (0,0)--(3,0);
      \draw[very thick,navy,->] (0,0)--(0,1);
    \end{scope}
    \node[navy,font=\small,anchor=south west] at (2.2,2.2) {$3\,\ket+$};
    \node[navy,font=\small,anchor=south east] at (-0.75,0.9) {$1\,\ket-$};
    \node[ink,font=\small,fill=white,inner sep=1.5pt,anchor=north west] at (0.9,-0.25) {unit circle};
    \node[figacc,font=\small,fill=white,inner sep=1.5pt,anchor=east] at (-0.2,3.0) {image of the circle under $A$};
  \end{tikzpicture}
  \caption{The operator $A=\begin{pmatrix}2&1\\1&2\end{pmatrix}$ of Example~\ref{ex:eigen} maps the unit circle to an ellipse. Its axes point along the orthonormal eigenvectors $\ket+$ and $\ket-$, and its semi-axes are the eigenvalues, $3$ and $1$.}
  \label{fig:ellipse}
\end{figure}

\begin{example}\label{ex:diag}
  Diagonalise $A=\begin{pmatrix}2&1\\1&2\end{pmatrix}$, and show that $\sx$ is, in the basis $\{\ket+,\ket-\}$, just the matrix $\sz$.
\end{example}
\begin{solution}
  \begin{steps}
    \item \emph{Eigenvectors.} By Example~\ref{ex:eigen}, the orthonormal eigenvectors are $\ket+$ (eigenvalue $3$) and $\ket-$ (eigenvalue $1$).
    \item \emph{The unitary.} Taking them as columns, $U=\tfrac1{\sqrt2}\begin{pmatrix}1&1\\1&-1\end{pmatrix}$.
    \item \emph{Diagonalise.} $U^\dagger AU=\begin{pmatrix}3&0\\0&1\end{pmatrix}$, so $A=3\ket+\!\bra++1\ket-\!\bra-$.
    \item \emph{The Pauli matrix.} Likewise $\sx\ket\pm=\pm\ket\pm$, so in the basis $\{\ket+,\ket-\}$ the matrix of $\sx$ is $\diag(1,-1)=\sz$.
    \item \emph{Interpretation.} This explains the identity $H\sx H=\sz$ of Example~\ref{ex:hadamard}: the matrix $U$ above is the Hadamard gate, which changes from the computational basis to $\{\ket\pm\}$, and $\sx=H\sz H$. The Pauli matrix $\sx$ that swaps $\ket0$ and $\ket1$ is, along $\ket\pm$, just ``multiply by $\pm1$''.
  \end{steps}
\end{solution}

\begin{lookahead}
  The spectral theorem is the mathematical content of a measurement. The eigenvalues of an observable are its possible outcomes, the eigenvectors are the states in which the outcome is certain, and the squared overlaps of a state with the eigenvectors will be the outcome probabilities (Postulate~3). The theorem also tells us that two observables that commute can be measured together, since they share an eigenbasis (Exercise~\ref{ex:commuting-basis}).
\end{lookahead}

\subsection{Projectors}

Spectral decompositions are made of pieces $\ket v\!\bra v$ which are orthogonal projections. It pays to study these on their own.

\begin{definition}[Orthogonal projector]\label{def:projector}
  An operator $P$ is an \term{orthogonal projector} if $P^2=P$ and $P^\dagger=P$.
\end{definition}

Let $W=\Ran P$. Then every vector splits uniquely as $\ket\psi=P\ket\psi+(\Id-P)\ket\psi$ with $P\ket\psi\in W$ and $(\Id-P)\ket\psi\in W^\perp$, where $W^\perp=\{\ket v:\braket wv=0\ \forall\,w\in W\}$ is the \term{orthogonal complement}, so that $\Hil=W\oplus W^\perp$ (Figure~\ref{fig:proj}). In summary:

\begin{proposition}\label{prop:proj}
  Let $P$ be an orthogonal projector with range $W$.
  \begin{enumerate}
    \item[(a)] $\Id-P$ is the orthogonal projector onto $W^\perp$;
    \item[(b)] the eigenvalues of $P$ are $0$ and $1$, with eigenspaces $W^\perp$ and $W$, and $\Tr P=\dim W$;
    \item[(c)] if $\{\ket{e_1},\dots,\ket{e_k}\}$ is an ONB of $W$, then $P=\sum_{i=1}^k\ket{e_i}\!\bra{e_i}$; in particular $\ket\psi\!\bra\psi$ is the projector onto $\ket\psi$ when $\norm\psi=1$;
    \item[(d)] \emph{best approximation:} $\norm{\psi-P\psi}\le\norm{\psi-w}$ for every $\ket w\in W$.
  \end{enumerate}
\end{proposition}
\begin{proof}
  First two facts. $P$ fixes every vector of $W$: if $\ket w=P\ket u$ then $P\ket w=P^2\ket u=\ket w$. And $\Ker P=W^\perp$: if $P\ket v=0$ and $\ket w=P\ket u\in W$, then $\braket wv=\braket{Pu}v=\braket u{Pv}=0$; conversely, if $\ket v\in W^\perp$ then $\norm{Pv}^2=\braket{Pv}{Pv}=\braket v{P^2v}=\braket v{Pv}=0$ because $P\ket v\in W$.
  \begin{steps}
    \item (a) $(\Id-P)^2=\Id-2P+P^2=\Id-P$ and $(\Id-P)^\dagger=\Id-P$, so $\Id-P$ is an orthogonal projector. Its range is $\Ker P=W^\perp$: every $(\Id-P)\ket u$ is killed by $P$, and every $\ket v\in\Ker P$ equals $(\Id-P)\ket v$.
    \item (b) If $P\ket v=\lambda\ket v$, then $\lambda\ket v=P^2\ket v=\lambda^2\ket v$, so $\lambda\in\{0,1\}$; the eigenvectors with $\lambda=1$ are the vectors fixed by $P$, which fill $W$, and those with $\lambda=0$ fill $\Ker P=W^\perp$. The trace is the sum of eigenvalues, that is, $\dim W$.
    \item (c) Let $Q=\sum_i\ket{e_i}\!\bra{e_i}$. For $\ket w\in W$ the expansion of Theorem~\ref{thm:parseval} gives $Q\ket w=\ket w$, and $Q\ket v=0$ for $\ket v\in W^\perp$; thus $Q$ and $P$ agree on $W$ and on $W^\perp$, hence everywhere.
    \item (d) Since $\psi-P\psi\in W^\perp$ and $P\psi-w\in W$ are orthogonal, Pythagoras gives $\norm{\psi-w}^2=\norm{\psi-P\psi}^2+\norm{P\psi-w}^2\ge\norm{\psi-P\psi}^2$. \qedhere
  \end{steps}
\end{proof}

\begin{figure}[htb]
  \centering
  \begin{tikzpicture}[>=stealth,scale=1.1]
    \draw[->,ink!70] (-0.5,0) -- (5,0) node[right] {$W$};
    \draw[->,ink!70] (0,-0.5) -- (0,3) node[above] {$W^\perp$};
    \draw[very thick,navy,->] (0,0) -- (3.4,2.1) node[above right] {$\ket\psi$};
    \draw[very thick,figacc,->] (0,0) -- (3.4,0) node[below right] {$P\ket\psi$};
    \draw[very thick,ink,->] (0,0) -- (0,2.1) node[left] {$(\Id-P)\ket\psi$};
    \draw[dashed,ink!60] (3.4,0) -- (3.4,2.1) -- (0,2.1);
  \end{tikzpicture}
  \caption{An orthogonal projector splits every vector into a part in $W$ and a part in $W^\perp$. $P\ket\psi$ is the point of $W$ closest to $\ket\psi$.}
  \label{fig:proj}
\end{figure}

By part (c), every projector is a sum of the rank-one pieces of a spectral decomposition. Grouping equal eigenvalues, the spectral theorem takes its most useful form:
\[
  H=\sum_\lambda\lambda\,P_\lambda,\qquad P_\lambda P_{\lambda'}=0\ \ (\lambda\neq\lambda'),\qquad\sum_\lambda P_\lambda=\Id,
\]
where $P_\lambda$ is the projector onto the eigenspace of $\lambda$. \emph{Every observable is a weighted sum of yes/no questions.}

\begin{xref}
  This is exactly how measurement will work in Postulate~3 (Section~\ref{sec:measurement}): measuring the observable $H=\sum_\lambda\lambda P_\lambda$ in the state $\ket\psi$ gives the outcome $\lambda$ with probability $\bra\psi P_\lambda\ket\psi=\norm{P_\lambda\psi}^2$, and afterwards the state is $P_\lambda\ket\psi/\norm{P_\lambda\psi}$.
\end{xref}

\begin{example}[A rank-one projector in the plane]\label{ex:proj-line}
  Find the matrix of the orthogonal projector onto the line spanned by $(1,2)$ in $\R^2$, and project $v=(3,1)$.
\end{example}
\begin{solution}
  \begin{steps}
    \item \emph{Normalise the direction.} $e=\frac1{\sqrt5}(1,2)$.
    \item \emph{The projector} (Proposition~\ref{prop:proj}(c)): $P=\ket e\!\bra e=\tfrac15\begin{pmatrix}1&2\\2&4\end{pmatrix}$.
    \item \emph{Check.} $P^2=\tfrac1{25}\begin{pmatrix}5&10\\10&20\end{pmatrix}=P$ and $\Tr P=1$.
    \item \emph{Project.} $Pv=\tfrac15\begin{pmatrix}1&2\\2&4\end{pmatrix}\begin{pmatrix}3\\1\end{pmatrix}=\tfrac15\begin{pmatrix}5\\10\end{pmatrix}=(1,2)$.
    \item \emph{The remainder} $(\Id-P)v=(2,-1)$ is orthogonal to $(1,2)$.
  \end{steps}
\end{solution}

\begin{example}\label{ex:proj-plane}
  Project $v=(1,2,3)$ orthogonally onto the plane $W=\{x+y+z=0\}\subset\R^3$.
\end{example}
\begin{solution}
  \begin{steps}
    \item \emph{The complement.} $W^\perp$ is the line spanned by the normal $n=\tfrac1{\sqrt3}(1,1,1)$, so the projector onto $W$ is $P=\Id-\ket n\!\bra n$.
    \item \emph{The component along $n$:} $\langle n,v\rangle n=\tfrac6{\sqrt3}\cdot\tfrac1{\sqrt3}(1,1,1)=(2,2,2)$.
    \item \emph{Subtract.} $Pv=(1,2,3)-(2,2,2)=(-1,0,1)$.
    \item \emph{Check.} It satisfies $x+y+z=0$.
  \end{steps}
\end{solution}

\begin{example}\label{ex:measure-preview}
  Let $\ket\psi=\frac1{\sqrt3}\ket0+\sqrt{\tfrac23}\ket1$. Write the spectral decomposition of $\sx$ in terms of projectors and compute $\bra\psi P_\pm\ket\psi$ for $P_\pm=\ket\pm\!\bra\pm$.
\end{example}
\begin{solution}
  \begin{steps}
    \item \emph{Decomposition.} From Example~\ref{ex:eigen}, $\sx=P_+-P_-$ with $P_\pm=\ket\pm\!\bra\pm=\tfrac12\begin{pmatrix}1&\pm1\\\pm1&1\end{pmatrix}$.
    \item \emph{Overlaps.} $\braket\pm\psi=\tfrac1{\sqrt2}\bigl(\tfrac1{\sqrt3}\pm\sqrt{\tfrac23}\bigr)=\tfrac{1\pm\sqrt2}{\sqrt6}$.
    \item \emph{Squared moduli.} $\bra\psi P_+\ket\psi=\frac{(1+\sqrt2)^2}{6}=\frac{3+2\sqrt2}{6}\approx0.971$ and $\bra\psi P_-\ket\psi=\frac{3-2\sqrt2}{6}\approx0.029$.
    \item \emph{Check.} They add up to $1$, as Parseval (and the completeness relation $P_++P_-=\Id$) demand.
  \end{steps}
\end{solution}

\subsection{Functions of operators}

Once an operator is diagonal, applying a function to it is easy: apply the function to each eigenvalue. We take this as the definition.

\begin{definition}[Functional calculus]\label{def:functional}
  Let $A$ be a Hermitian or unitary (more generally normal) operator, $A=\sum_i\lambda_i\ket{v_i}\!\bra{v_i}$, and let $f$ be a function defined on its spectrum. Then
  \[
    f(A)=\sum_if(\lambda_i)\ket{v_i}\!\bra{v_i}.
  \]
\end{definition}

If $f(z)=\sum_nc_nz^n$ is a power series, this agrees with $\sum_nc_nA^n$, because $A^n=\sum_i\lambda_i^n\ket{v_i}\!\bra{v_i}$. In short: \emph{apply $f$ to the eigenvalues, keep the eigenvectors.} The operator $f(A)$ has the same eigenvectors as $A$, $(fg)(A)=f(A)g(A)$, and $f(A)$ is Hermitian when $f$ is real-valued on the spectrum. Two cases matter most. The first is the exponential of $\ii$ times a Hermitian operator.

\begin{proposition}\label{prop:expunitary}
  If $K$ is Hermitian, then $e^{\ii K}$ is unitary. Conversely, every unitary $U$ can be written $U=e^{\ii K}$ with $K$ Hermitian.
\end{proposition}
\begin{proof}
  \begin{steps}
    \item \emph{Diagonalise.} Let $K=\sum_j\kappa_jP_j$ with $\kappa_j\in\R$ and $P_j$ orthogonal projectors; then $e^{\ii K}=\sum_je^{\ii\kappa_j}P_j$.
    \item \emph{Unitarity.} $\bigl(e^{\ii K}\bigr)^\dagger e^{\ii K}=\sum_je^{-\ii\kappa_j}e^{\ii\kappa_j}P_j=\sum_jP_j=\Id$.
    \item \emph{The converse.} If $U=\sum_je^{\ii\theta_j}P_j$ (Theorem~\ref{thm:spectrum}(b)), take $K=\sum_j\theta_jP_j$. \qedhere
  \end{steps}
\end{proof}

The second case is a projector, for which the spectrum is $\{0,1\}$: $f(P)=f(0)(\Id-P)+f(1)P$.

\begin{example}\label{ex:funcs}
  (a) Compute $\sqrt A$ for $A=\begin{pmatrix}2&1\\1&2\end{pmatrix}$. (b) Show that $e^{-\ii\theta\sx/2}=\cos\frac\theta2\,\Id-\ii\sin\frac\theta2\,\sx$.
\end{example}
\begin{solution}
  \begin{steps}
    \item (a) \emph{Spectral decomposition.} From Example~\ref{ex:diag}, $A=3P_++1P_-$ with $P_\pm=\tfrac12\begin{pmatrix}1&\pm1\\\pm1&1\end{pmatrix}$.
    \item (a) \emph{Apply $f(x)=\sqrt x$:} $\sqrt A=\sqrt3\,P_++P_-=\frac12\begin{pmatrix}\sqrt3+1&\sqrt3-1\\\sqrt3-1&\sqrt3+1\end{pmatrix}$.
    \item (a) \emph{Check.} The diagonal entry of $(\sqrt A)^2$ is $\tfrac14\bigl((\sqrt3+1)^2+(\sqrt3-1)^2\bigr)=2$ and the off-diagonal one is $\tfrac24\cdot2=1$.
    \item (b) \emph{Spectral decomposition of $\sx$:} $\sx=P_+-P_-$ and $P_++P_-=\Id$.
    \item (b) \emph{Apply the exponential:} $e^{-\ii\theta\sx/2}=e^{-\ii\theta/2}P_++e^{\ii\theta/2}P_-=\cos\tfrac\theta2(P_++P_-)-\ii\sin\tfrac\theta2(P_+-P_-)=\cos\tfrac\theta2\,\Id-\ii\sin\tfrac\theta2\,\sx$.
    \item (b) \emph{Meaning.} This is a unitary operator by Proposition~\ref{prop:expunitary}, the rotation of a qubit by an angle $\theta$ about the $x$-axis. The same computation with $\sz$ in place of $\sx$ gives $e^{-\ii\theta\sz/2}=\diag(e^{-\ii\theta/2},e^{\ii\theta/2})$. Curiously, for $\theta=2\pi$ it equals $-\Id$: a full turn gives a factor $-1$, a global phase.
  \end{steps}
\end{solution}

\begin{example}[A phase gate as an exponential]\label{ex:phase-exp}
  Show that $S=\diag(1,\ii)$ equals $e^{\ii\pi/4}\,e^{-\ii\frac\pi4\sz}$ and find a Hermitian $K$ with $S=e^{\ii K}$.
\end{example}
\begin{solution}
  \begin{steps}
    \item \emph{Exponential of $\sz$.} $e^{-\ii\frac\pi4\sz}=\diag(e^{-\ii\pi/4},e^{\ii\pi/4})$.
    \item \emph{Multiply by $e^{\ii\pi/4}$:} $\diag(1,e^{\ii\pi/2})=\diag(1,\ii)=S$.
    \item \emph{A Hermitian logarithm.} $S=e^{\ii K}$ with $K=\diag(0,\tfrac\pi2)=\tfrac\pi4(\Id-\sz)$ (Proposition~\ref{prop:expunitary}; $K$ is determined up to adding multiples of $2\pi$ to the eigenvalues).
  \end{steps}
\end{solution}

\begin{pitfall}
  Functions of operators are \emph{not} applied entrywise: $e^A$ is not the matrix of the numbers $e^{A_{ij}}$. And the law $e^{A+B}=e^Ae^B$ holds only when $AB=BA$. Since $\sx\sz\neq\sz\sx$, in general $e^{\ii(\alpha\sx+\beta\sz)}\neq e^{\ii\alpha\sx}e^{\ii\beta\sz}$.
\end{pitfall}

A case that will dominate Chapter~\ref{ch:postulates} is the Hamiltonian $H$ of a system, a Hermitian operator whose eigenvalues $E_i$ are the allowed energies, $H=\sum_iE_i\ket{E_i}\!\bra{E_i}$. Its exponential
\[
  e^{-\ii Ht/\hbar}=\sum_ie^{-\ii E_it/\hbar}\ket{E_i}\!\bra{E_i}
\]
is unitary, and it multiplies each energy eigenstate $\ket{E_i}$ by the phase $e^{-\ii E_it/\hbar}$. This is how Schr\"odinger's equation will become a unitary map (Proposition~\ref{prop:schrodinger}), and why a system with a single energy does not change (Section~\ref{sec:levels}).

\begin{lookahead}
  The functional calculus has one more famous customer. For a density operator $\rho$ (Section~\ref{sec:density}), the operator $\rho\log\rho$ is defined by $f(x)=x\log x$, and $-\Tr(\rho\log\rho)$ is the \emph{von Neumann entropy} of Chapter~\ref{ch:info}, the quantum measure of how much we do not know about a state.
\end{lookahead}

\subsection{The picture so far}

Every concept of this chapter has a role in a quantum computer. The table collects them.

\begin{center}
\small
\begin{tabular}{@{}p{0.27\textwidth}p{0.33\textwidth}p{0.3\textwidth}@{}}
  \toprule
  Mathematics & What it means & In a quantum computer\\
  \midrule
  unit vector $\ket\psi\in\Hil$ & the state of the system & the data\\
  inner product $\braket\varphi\psi$ & an amplitude; squared, a probability & readout of results\\
  orthonormal basis $\{\ket{e_i}\}$ & distinguishable outcomes & a choice of ``question''\\
  Hermitian $H=H^\dagger$ & an observable; eigenvalues are the outcomes & measurement\\
  projector $P$ & the yes/no question ``is it in $W$?'' & measurement step\\
  unitary $U^\dagger U=\Id$ & reversible evolution & quantum gates\\
  $f(A)$, e.g.\ $e^{-\ii Ht/\hbar}$ & dynamics generated by $H$ & gates from physics\\
  \bottomrule
\end{tabular}
\end{center}

In one sentence: \emph{a quantum computer prepares a unit vector, rotates it with unitaries, and reads it with projectors} (Figure~\ref{fig:pipeline}). Chapter~\ref{ch:postulates} turns this sentence into precise postulates.

\begin{figure}[htb]
  \centering
  \begin{tikzpicture}[>=stealth,node distance=0pt,
      box/.style={draw=navy,thick,fill=navy!8,minimum width=1.2cm,minimum height=0.8cm,font=\small}]
    \node[font=\small] (in) at (0,0) {$\ket0$};
    \node[box] (u1) at (2,0) {$U_1$};
    \node[box] (u2) at (3.6,0) {$U_2$};
    \node[font=\small] (dots) at (5.1,0) {$\cdots$};
    \node[box] (uk) at (6.6,0) {$U_k$};
    \node[draw=figacc,thick,fill=figacc!10,minimum width=1.4cm,minimum height=0.8cm,font=\small] (m) at (8.5,0) {$\{P_\lambda\}$};
    \draw[->] (in)--(u1); \draw[->] (u1)--(u2); \draw[->] (u2)--(dots); \draw[->] (dots)--(uk); \draw[->] (uk)--(m);
    \node[font=\footnotesize,ink] at (0,-0.9) {prepare};
    \node[font=\footnotesize,ink] at (4.3,-0.9) {evolve (reversible, unitary)};
    \node[font=\footnotesize,ink] at (8.5,-0.9) {measure (irreversible)};
  \end{tikzpicture}
  \caption{The skeleton of every quantum algorithm: prepare a state, apply unitaries, measure with projectors. An algorithm is a choice of $U_1,\dots,U_k$ such that the right answer has probability close to $1$ at the end.}
  \label{fig:pipeline}
\end{figure}

\begin{exerciseset}
  \exercise{1} Find the adjoint of $\begin{pmatrix}1&\ii\\2&3-\ii\end{pmatrix}$ and of each Pauli matrix. Compute $M^t$, $\bar M$ and $M^\dagger$ for $M=\begin{pmatrix}1&2\ii\\3-\ii&4\end{pmatrix}$.

  \exercise{2} Prove the rules of the adjoint: $(A^\dagger)^\dagger=A$, $(\lambda A)^\dagger=\bar\lambda A^\dagger$, $(AB)^\dagger=B^\dagger A^\dagger$, $(A^{-1})^\dagger=(A^\dagger)^{-1}$; and in Dirac notation $(A\ket\psi)^\dagger=\bra\psi A^\dagger$.

  \exercise{2} Show that the diagonal entries of a Hermitian matrix are real, that $A^\dagger A$ and $A+A^\dagger$ and $\ii(A-A^\dagger)$ are Hermitian for every $A$, and that $\ii[A,B]$ is Hermitian when $A,B$ are.

  \exercise{2} Let $\{\ket{e_i}\}$ be a non-orthonormal basis of $\C^2$, say $\ket{e_1}=\ket0$, $\ket{e_2}=\ket+$. Compute the matrix of $A=\sx$ in this basis, and the matrix of $A^\dagger$ in this basis; show that the latter is not the conjugate transpose of the former.

  \exercise{2} Show that if $U$ is unitary then $U$ preserves inner products, $\abs{\det U}=1$, and that the columns of $U$ form an orthonormal basis. Show that the general $2\times2$ unitary is $e^{\ii\alpha}\begin{pmatrix}a&-\bar b\\b&\bar a\end{pmatrix}$ with $\abs a^2+\abs b^2=1$.

  \exercise{1} Find eigenvalues and eigenvectors over $\C$ of $\begin{pmatrix}2&1\\1&2\end{pmatrix}$, $\begin{pmatrix}1&\ii\\-\ii&1\end{pmatrix}$, $\begin{pmatrix}0&1\\0&0\end{pmatrix}$, $\begin{pmatrix}1&1\\0&1\end{pmatrix}$ and $\diag(2,2,3)$. Which have real eigenvalues? Which have an eigenbasis? An orthonormal one?

  \exercise{1} Find eigenvalues and normalised eigenvectors of $\sx,\sy,\sz$ and check that their eigenbases are $\{\ket\pm\}$, $\{\ket{\pm\ii}\}$ and $\{\ket0,\ket1\}$ up to phases.

  \exercise{2} For $A\in M_2(\C)$ show that $p(\lambda)=\lambda^2-\Tr(A)\lambda+\det A$, hence $\lambda_1+\lambda_2=\Tr A$ and $\lambda_1\lambda_2=\det A$. Generalise to $n\times n$ (with multiplicities).

  \exercise{2} Prove the following. (a) A Hermitian operator has real eigenvalues. (b) Eigenvectors of a Hermitian operator for different eigenvalues are orthogonal. (c) A unitary operator has eigenvalues of modulus $1$. (d) Eigenvectors of a unitary operator for distinct eigenvalues are orthogonal (hint: compute $\braket{Uv}{Uw}$ in two ways).

  \exercise{2} If $A\ket v=\lambda\ket v$, show that $A^k\ket v=\lambda^k\ket v$, $p(A)\ket v=p(\lambda)\ket v$ for every polynomial $p$, and $A^{-1}\ket v=\lambda^{-1}\ket v$ when $A$ is invertible. Find the eigenvalues of $A^{10}$ for $A=\begin{pmatrix}2&1\\1&2\end{pmatrix}$.

  \exercise{2} Show that eigenvectors of any operator for distinct eigenvalues are linearly independent. Deduce that an $n\times n$ matrix with $n$ distinct eigenvalues is diagonalisable.

  \exercise{2} Diagonalise $\begin{pmatrix}0&-\ii\\\ii&0\end{pmatrix}$ and $\begin{pmatrix}1&1+\ii\\1-\ii&2\end{pmatrix}$, finding a unitary $U$ with $U^\dagger AU$ diagonal.

  \exercise{2} Show that the rotation matrix $R_\theta$ ($\theta\notin\pi\Z$) has no real eigenvector but has complex eigenvalues $e^{\pm\ii\theta}$. Why does this make $\C$ the natural field for spectral theory?

  \exercise{3} An operator $A$ is \emph{positive} ($A\ge0$) if $\bra\psi A\ket\psi\ge0$ for all $\ket\psi$. Show that over $\C$ a positive operator is Hermitian; that a Hermitian $A$ is positive iff its eigenvalues are $\ge0$; and that $B^\dagger B\ge0$ for every $B$ and every projector is positive. (We return to positive operators in Section~\ref{sec:measurement}.)

  \exercise{3} Let $A,B$ be Hermitian with $[A,B]=0$ and suppose $A$ has distinct eigenvalues. Show that every eigenvector of $A$ is an eigenvector of $B$. Check on $A=\sz$ and $B=\Id+2\sz$, and explain why this fails for $A=\sz$, $B=\sx$.

  \exercise{1} For $\ket\psi=\ket{+\ii}$ and $\ket\varphi=\cos\frac\theta2\ket0+e^{\ii\phi}\sin\frac\theta2\ket1$, write $P=\ket\varphi\!\bra\varphi$ as a matrix and show $P=\tfrac12(\Id+n\cdot\sigma)$ with $n=(\sin\theta\cos\phi,\sin\theta\sin\phi,\cos\theta)$.

  \exercise{2} For a unit vector $\ket\psi$ and $P=\ket\psi\!\bra\psi$ prove $P=P^\dagger$, $P^2=P$, $\Tr P=1$, find its eigenvalues and eigenspaces, and show $\Tr(P\ket\varphi\!\bra\varphi)=\abs{\braket\varphi\psi}^2$.

  \exercise{2} Find the matrix of the orthogonal projector onto the line spanned by $(1,-1,2)$ in $\R^3$, and project $(2,0,1)$ onto it.

  \exercise{2} For a subspace $W<\Hil$ prove: $W^\perp$ is a subspace, $W\cap W^\perp=\{0\}$, $\dim W+\dim W^\perp=\dim\Hil$, $(W^\perp)^\perp=W$. Compute $W^\perp$ for $W=\Span\{\ket0+\ii\ket1\}\subset\C^2$ and for $W=\Span\{(1,\ii,0),(0,1,1)\}\subset\C^3$.

  \exercise{2} Let $\{\ket{e_1},\dots,\ket{e_k}\}$ be an ONB of $W$ and $P=\sum_i\ket{e_i}\!\bra{e_i}$. Show $P^2=P=P^\dagger$, $\Ran P=W$, $\Ker P=W^\perp$ and $\Tr P=\dim W$; show that $P$ does not depend on the chosen ONB of $W$; and show that every orthogonal projector is of this form.

  \exercise{2} Project $f(x)=x^2$ onto $\Span\{1,x\}$ in $L^2[-1,1]$, and $(1,2,3)$ onto $x+y+z=0$ in $\R^3$. In each case verify Proposition~\ref{prop:proj}(d).

  \exercise{2} For orthogonal projectors $P_1,P_2$ show: $P_1+P_2$ is a projector iff $P_1P_2=0$; $P_1P_2$ is a projector iff $P_1P_2=P_2P_1$; $\Id-P_1$ is the projector onto $(\Ran P_1)^\perp$.

  \exercise{2} Measuring $A=\sum_\lambda\lambda P_\lambda$ on $\ket\psi$ gives $\lambda$ with probability $p(\lambda)=\bra\psi P_\lambda\ket\psi$ and leaves the state $P_\lambda\ket\psi/\sqrt{p(\lambda)}$. Show $\sum_\lambda p(\lambda)=1$, and compute the probabilities and post-measurement states for $A=\sx$ and $\ket\psi=\frac1{\sqrt3}\ket0+\sqrt{\tfrac23}\ket1$.

  \exercise{2} Find $\sqrt{\begin{pmatrix}2&1\\1&2\end{pmatrix}}$ and $e^{\ii\pi\sx/2}$. Show that for $P$ a projector, $f(P)=f(0)(\Id-P)+f(1)P$, and compute $e^{\ii\theta P}$.

  \exercise{2} Show that $e^{-\ii\theta\,n\cdot\sigma/2}=\cos\frac\theta2\Id-\ii\sin\frac\theta2\,n\cdot\sigma$ for a unit vector $n$. Compute the rotations $R_x(\pi)$, $R_y(\pi/2)$ and $R_z(2\pi)$ and show that $R_z(2\pi)=-\Id$.

  \exercise{2} Show that every unitary $2\times2$ matrix with determinant $1$ is of the form $\cos\frac\theta2\Id-\ii\sin\frac\theta2\,n\cdot\sigma$, and find $\theta,n$ for $-\ii\sx$ and for $e^{-\ii\pi/4}\diag(1,\ii)$.

  \exercise{2} Write the phase gate $T=\diag(1,e^{\ii\pi/4})$ as $e^{\ii\gamma}e^{-\ii\theta\sz/2}$, finding $\gamma$ and $\theta$.

  \exercise{3} Show that $e^{\ii(\alpha\sx+\beta\sz)}\neq e^{\ii\alpha\sx}e^{\ii\beta\sz}$ for generic $\alpha,\beta$, by comparing the traces or the $(1,1)$ entries at second order in $\alpha,\beta$. When is equality true?
\end{exerciseset}

\section*{Further reading}

\begin{itemize}
  \item Axler~\cite{axler2015}: the cleanest modern account of finite-dimensional linear algebra, with the same emphasis on operators and inner products as this chapter.
  \item Horn and Johnson~\cite{horn2013}: the reference for matrix analysis, including the singular value decomposition used in Chapter~\ref{ch:postulates}.
  \item Griffiths and Schroeter~\cite{griffiths2018}: the physics behind Section~\ref{sec:motivation}, at an accessible level.
  \item Nielsen and Chuang~\cite{nielsenchuang2010}, Chapter~2: the standard treatment of the linear algebra of quantum computing; compare its notation with ours.
\end{itemize}

